\documentclass[11pt]{article}
\ifdefined\pdftrailerid\pdftrailerid{}\fi
\usepackage[T1]{fontenc}
\usepackage[utf8]{inputenc}
\usepackage{lmodern}
\usepackage[margin=1in]{geometry}
\usepackage{amsmath,amssymb,amsthm,mathtools}
\usepackage{enumitem,booktabs,array,longtable}
\usepackage{xcolor,microtype,etoolbox,xurl}
\usepackage[colorlinks=true,linkcolor=blue!45!black,citecolor=green!35!black,urlcolor=blue!55!black]{hyperref}
\usepackage[capitalize,nameinlink]{cleveref}
\hypersetup{pdftitle={Source Covers and Barriers to Tape Transposition},pdfauthor={Hans Brende and GPT-6 Astra (OpenAI)},pdfsubject={Complete source covers and restricted nonlinear transposition impossibility},pdfkeywords={matrix transposition, multitape Turing machines, source covers, nonlinear coding}}
\numberwithin{equation}{section}
% ---------------------------------------------------------------------------
% Theorem environments
% ---------------------------------------------------------------------------
\theoremstyle{plain}
\newtheorem{theorem}{Theorem}[section]
\newtheorem{lemma}[theorem]{Lemma}
\newtheorem{corollary}[theorem]{Corollary}
\newtheorem{proposition}[theorem]{Proposition}
\newtheorem{conjecture}[theorem]{Conjecture}
\newtheorem{question}[theorem]{Question}
\newtheorem{problem}[theorem]{Problem}
\theoremstyle{definition}
\newtheorem{definition}[theorem]{Definition}
\newtheorem{remark}[theorem]{Remark}
\newtheorem{example}[theorem]{Example}

% Formalized results use outlines; unformalized results retain complete written proofs.
\newenvironment{outline}[1][Outline of the formal proof]{\begin{proof}[#1]}{\end{proof}}
% Results that are not formalized get complete proofs, labelled as such.
\newenvironment{informalproof}{\begin{proof}[Proof (not formalized)]}{\end{proof}}

% ---------------------------------------------------------------------------
% Lean references
%   \lean{name}            a declaration cited with its source paper or namespace
%   \formal{file}{names}   a relative source file and comma-separated declaration names
%   \partformal{file}{names}{gap}
%                          a numbered result whose formal counterpart covers all but the
%                          part described in {gap}; the gap is proved or cited in the text
%   \notformalized         a numbered result that has no formal counterpart
%   \cited{where}          a result of an earlier paper used without proof: where it is proved,
%                          and whether it is Lean-verified there
% ---------------------------------------------------------------------------
% Lean names are typeset like paths, with line breaks allowed after "_", "." and "/".
\DeclareUrlCommand\leanname{\def\UrlFont{\ttfamily\upshape}}
\renewcommand{\UrlBreaks}{\do\_\do\.\do\/\do\-}
\newcommand{\lean}[1]{\leanname{#1}}
\newcommand{\leanfile}[1]{\leanname{#1}}
\newcommand{\leanlist}[1]{\begingroup\def\leansep{}%
  \renewcommand*{\do}[1]{\leansep\leanname{##1}\def\leansep{, }}\docsvlist{#1}\endgroup}
\newcommand{\formal}[2]{\par\noindent{\small\upshape\raggedright\textsf{Lean:} \leanlist{#2}\unskip
  \nobreak\hfil\penalty50\hskip1em\hbox{}\nobreak\hfill[\leanfile{#1}]\parfillskip=0pt\par}%
  \smallskip\noindent\ignorespaces}
\newcommand{\partformal}[3]{\par\noindent{\small\upshape\raggedright\textsf{Lean (partial):} \leanlist{#2}\unskip
  \nobreak\hfil\penalty50\hskip1em\hbox{}\nobreak\hfill[\leanfile{#1}]\parfillskip=0pt\par}%
  \par\noindent{\small\upshape\textsf{Not formalized:} #3}\par\smallskip\noindent\ignorespaces}
\newcommand{\notformalized}{\par\noindent{\small\upshape\textsf{Not formalized;} a complete proof is given
  below.}\par\smallskip\noindent\ignorespaces}
\newcommand{\cited}[1]{\par\noindent{\small\upshape\textsf{Cited:} #1}\par\smallskip\noindent\ignorespaces}

% ---------------------------------------------------------------------------
% Notation
% ---------------------------------------------------------------------------
\newcommand{\N}{\mathbb{N}}
\newcommand{\Z}{\mathbb{Z}}
\newcommand{\Q}{\mathbb{Q}}
\newcommand{\R}{\mathbb{R}}
\newcommand{\C}{\mathbb{C}}
\newcommand{\F}{\mathbb{F}}
\newcommand{\T}{\mathsf{T}}
\newcommand{\clog}[1]{\lceil \log_2 #1 \rceil}
\newcommand{\abs}[1]{\lvert #1 \rvert}
\newcommand{\card}[1]{\lvert #1 \rvert}
\newcommand{\eps}{\varepsilon}
\newcommand{\prev}{\mathrm{prev}}
\newcommand{\cost}{\mathrm{cost}}
\newcommand{\gn}{\mathrm{gn}}
\newcommand{\Sh}{\mathrm{Sh}}
\newcommand{\cl}{\mathrm{cl}}
\newcommand{\Aut}{\mathrm{Aut}}
\newcommand{\minrk}{\mathrm{minrk}}
\newcommand{\MAIS}{\mathrm{MAIS}}
\newcommand{\Crux}{\mathrm{Crux}}
\newcommand{\TCP}{\mathrm{TCP}}
\DeclareMathOperator{\rank}{rank}
\DeclareMathOperator{\dist}{dist}
\DeclareMathOperator{\poly}{poly}

% Notation specific to this written-proof manuscript.
\DeclareMathOperator{\Var}{Var}
\newcommand{\gaparea}{\mathcal A}

\setlist{itemsep=3pt,topsep=5pt}
\setlength{\emergencystretch}{2em}
\title{Source Covers and Barriers to Tape Transposition}
\author{Hans Brende \and GPT-6 Astra (OpenAI)}
\date{October 2026}
\begin{document}
\maketitle
\begin{abstract}
We study linear-time transposition of a $d\times d$ matrix with
$b(d)=\max(1,\lceil\log_2d\rceil)$-bit entries on a fixed finite number of tapes.
A source certificate must retain every input on which an output can depend
and charge the information used to simplify that dependence.
We prove a finite correctness bound using the circular gaps in complete raw-source sets,
then strengthen it by fractional covers, anchored approximation, and shared raw revelations.
For fixed linear runtime and capacity bounds and a fixed bound on the
number of disjoint-span classes in the full family system, we obtain a
uniform positive correctness deficit in large dimensions for a two-class
core plus at most one arbitrary exceptional family and sufficiently small
remaining families. No size envelope on the core is needed.
With a fixed number $q$ of exceptional families, near-full correctness forces a
quantitative lower bound on the tail of the sorted core-family sizes.
We also give uniform exclusions under economical source-cover or global-sharing conditions.
With two exceptions, near-full correctness forces the dimension lower limit
of the global continuation count to grow faster than linearly in the prefix
cutoff, at every sufficiently large fixed protected scale.
A native example shows that quadratic pointwise cover growth is possible even for
two serial exceptional families. Two complementary results identify the movement cost
of a specified coded-caching query interface and the information needed to isolate
narrow consecutive column groups causally.
Under the explicit block-capacity bound, fixed causally autonomous groups
of size $o(\log d)$ require helper entropy plus correctness deficit at least $L-o(L)$.
All 32 numbered results and 6 model definitions are formalized in Lean.
The general transposition question remains open: the sufficient certificate conditions
are not proved for arbitrary tape computations.

\end{abstract}
\begin{quote}\small
\textbf{Contributions.} Hans Brende chose the problem and directed the research.
GPT-6 Astra (OpenAI) developed the mathematical arguments, Lean formalization
and text under his direction. This draft is partially formalized; each result
states its proof status.
\end{quote}
\clearpage
\begingroup\small\tableofcontents\endgroup
\clearpage
\section{Introduction}\label{sec:introduction}

Let $b(d)=\max(1,\lceil\log_2d\rceil)$ and $L=d^2b(d)$.
Question (T) asks whether one deterministic multitape Turing machine, with a fixed
finite number of tapes and a fixed finite alphabet, transposes all such matrices
in $O(L)$ time. The input and output are written in row-major order, so the
output lists the original columns consecutively. The question is motivated in
part by the reductions between transposition and integer multiplication of
Harvey and van der Hoeven~\cite{HvdH2025}.

An obstruction to routing alone does not settle (T). A tape computation may
store nonlinear functions of many input bits and later combine several such
records. The first paper of this series isolates that distinction
\cite{FirstPaper2026}. The depth-reduction program of the fourth and eighth
papers cuts a dependency graph by revealing a small transcript, and then asks
for a coding bound on the remaining recovery relation
\cite{FourthPaper2026,EighthPaper2026}. The seventh paper proves half-rate
bounds under explicit elimination hypotheses~\cite{SeventhPaper2026}.
This paper instead constructs complete source certificates for particular
return geometries and proves that their information deficit exceeds the
cost of the revelations.

\subsection{The argument and its scope}

Fix one full computation geometry and a nonempty set $\mathcal C$ of inputs
correctly transposed on that geometry. Write
\[
 \Delta=L-\log_2|\mathcal C|.
\]
A transcript of at most $g$ bits partitions $\mathcal C$ into at most $2^g$
fibres. If, on every fibre, the complementary coordinates determine a set
$A$ of input coordinates, then
\begin{equation}\label{eq:intro-paid-deficit}
 \Delta\ge |A|-g.
\end{equation}
Our work is to supply $A$ and $g$ from the same actual geometry. Complete
source sets retain direct input reads, state and work-tape parents, and
unpaid parents entering an output vertex whose record has been revealed.
The certificate permits nonlinear computations and correlated correctness
sets. It never assigns independent information to different source paths.

The main geometric hypothesis concerns complete return families. A
$\mathrm{Q}_2$ core has two classes of pairwise disjoint full return spans;
the definition and the strict clipping convention are given in
\cref{sec:source-engine}. An exceptional family is one complete family,
not one work tape. In addition to this decomposition we retain fixed
runtime and whole-step capacity constants and a fixed bound on the number
of disjoint-span classes in the full family system.

\begin{enumerate}
\item \emph{A finite source certificate.} A single cyclic ordering of the raw
coordinates produces an acyclic recoverable set. The size of that set is
controlled by the circular gap areas of complete source banks, with explicit
movement, grid and output-stripe errors. Fractional covers and anchored
approximation improve its geometric term without changing the recovery bank.
\item \emph{Uniform restricted impossibility.} A $\mathrm{Q}_2$ core plus one
arbitrary exceptional family cannot have both $\Delta=o(L)$ and maximum
outside-family size $o(L)$. For $q$ fixed exceptions, every limiting sorted
hop profile has tail at least a positive multiple of $R^{1-q}$.
Consequently a common size envelope $H/i^p$ with $p>q$ is excluded.
\item \emph{Further sufficient conditions.} Economical fractional source
covers, or sufficiently shared complete raw support paid for once across
all outputs, give additional uniform correctness-deficit theorems.
The quantifiers keep every finite cut parameter fixed before dimension
increases. These are criteria on a native class, not assertions that an
arbitrary machine satisfies them.
\item \emph{Construction restrictions and limits.} We exhibit a quadratic
pointwise cover obstruction, prove a movement bound for a specified query
interface even after reordering, and show that isolating fixed groups of
$o(\log d)$ consecutive columns causally requires an almost complete helper.
Each result states the interface it constrains.
\end{enumerate}

For an oblivious transposer correct on every input, $\Delta=0$; a positive
uniform deficit therefore excludes its geometry. For an adaptive machine,
passing to one geometry may already discard much of the input entropy.
We do not infer a large common-geometry correctness set without paying
that loss. Nor do we prove that the stronger cover criteria exclude a
native class outside the union of all earlier exclusions.

\subsection{Relation to the earlier papers and other work}

The source-bank construction uses actual return families and the owner
ordering studied in the eighth paper~\cite{EighthPaper2026}, and the
dependency semantics of the fourth~\cite{FourthPaper2026}. We give the
additional packing, residual-cut and raw-source proofs here.
The fifth paper develops complete recovery relations and geometric
restrictions on their codes~\cite{DirectedRecovery2026}. Its distinction
between a complete relation and a convenient list of selected paths is
essential throughout the present argument. The local coding barriers of
the second paper~\cite{SecondPaper2026} concern different functional
interfaces; they do not supply the source covers assumed here.

The fixed-geometry deficit theorems are complementary to the depth
reductions of the eighth paper. A depth bound alone leaves a coding
question. Under our different geometric hypotheses we complete the
deficit calculation. This is not a general improvement of the seventh
paper's half-rate theorem or a proof of the unrestricted coding crux.

The movement example uses an established placement delivery array from
coded caching~\cite{YanChengTangChen2017}; we credit that code and prove
only the stated additional query-movement assertion. The causal-group
argument uses finite max-flow/min-cut~\cite{FordFulkerson1956} and the
first paper's native routing metric. Its additional content is a paid
extraction of partial routing from a locally autonomous computation on
correlated helper fibres. We make no claim that the elementary covering,
entropy or flow principles are new.

\subsection{Proof status and organization}

All 32 numbered results and 6 numbered definitions
are formalized in Lean. They reuse the eighth paper's family geometry,
shared incidence, interval-cover, interval-graph, finite-orbit,
shifted-grid, causal-recovery, circular-gap and stripe-counting lemmas,
the second paper's finite-entropy library, the first paper's native
counting library, and Mathlib's compactness, summability and bounded-limit results. Their Lean statements are
authoritative; the accompanying outlines are guides to those proofs.
Cited results state their provenance and formal status
where used. The index in \cref{sec:theorem-index} identifies every exact statement,
its current status and the verified declarations. The index itself is not a
verification certificate.

\Cref{sec:source-engine} establishes the model, source bank and finite
freezing bound. \Cref{sec:source-covers} supplies the fractional and approximate
cover theorem. \Cref{sec:uniform} derives the uniform profile consequences;
\cref{sec:shared-support} treats raw support shared across outputs.
\Cref{sec:limitations,sec:movement,sec:causal-groups} explain the limits and
additional construction restrictions. \Cref{sec:roadmap} identifies the
remaining native geometry and computation questions.

\section{Complete source banks and paid recovery}\label{sec:source-engine}

This section proves the common finite statement used in the later exclusions.
The bank must cover every surviving raw source. A list of selected return
endpoints, without the state and unselected work dependencies, does not suffice.

\subsection{Fixed geometries and complete return families}

Fix integers $d,b\ge1$, and put $L=d^2b$, $W=db$ and
$\Omega=\{0,\ldots,L-1\}$. The finite results allow arbitrary $b$; the
application to~(T) takes $b=b(d)=\max(1,\lceil\log_2d\rceil)$. A
\emph{fixed complete geometry} consists of the input-head and work-head
walks on $\{0,\ldots,T-1\}$, the halting time, and the distinct increasing
append times $t_0<\cdots<t_{L-1}$, where $T\le cL$. The walks move by at
most one cell per step. For the Lean model, extend each walk constantly
after $T-1$; this preserves its unit-step property and all hops before $T$.
The output at time $t_r$ demands input coordinate
\begin{equation}\label{eq:transpose-demand}
 \pi(jW+ib+k)=iW+jb+k
 \qquad(0\le i,j<d,\ 0\le k<b).
\end{equation}
The geometry can be an oblivious schedule or one trace shared by a set
of inputs. We make no assertion that an arbitrary adaptive transposer has
a trace containing nearly all its input information.

The dependency graph has the state edges $t\to t-1$ and every work-head
previous-visit edge $u\to p$, where $p$ is the last earlier visit of that
head to its cell at time~$u$. A work edge of length one coincides with a
state edge. Coincident edges may be identified for reachability, but an
edge is retained whenever any physical dependency supplies it. Input
reads are sources, not previous-visit edges. The complete value at a time
contains its state after the step, its written work symbols and any output
bit. Each such value has at most $\beta$ bits of capacity, for a fixed
$\beta>0$; this includes all data needed by later steps.

A \emph{genie} is a set $\Gamma$ of time vertices and its transcript
$G_\Gamma$ consists of those complete values. Thus
$H(G_\Gamma)\le\beta\card\Gamma$. A backward path
$v_0\to\cdots\to v_s$ survives the genie if
$v_1,\ldots,v_s\notin\Gamma$. In particular, its root is retained even
when $v_0\in\Gamma$: we retain its direct read and every incoming
dependency whose lower endpoint is unpaid. The full source relation at
an output contains every input coordinate read at the end of such a path,
including a path of length zero. The usual recursive evaluation of the
dependency graph shows that this full source relation and $G_\Gamma$
determine the output. This is the cone convention of
\cite[Section~4.1]{FourthPaper2026}.

Let $\mathcal C\subseteq\{0,1\}^{\Omega}$ be a nonempty set correctly
transposed on the fixed geometry, and write
\begin{equation}\label{eq:correctness-deficit}
 \Delta=L-\log_2\card{\mathcal C}.
\end{equation}
All selections of families, time maps and paid vertices below depend on
the fixed geometry, not on the input values within $\mathcal C$. We use
ordinary base-two Shannon entropy only for finite random variables.

Write $x(t)$ for the input-head position clamped to $[0,L-1]$. This does
not increase its total variation. A read outside the input is a known
blank; including the bit at the clamped position supplies extra
information and is therefore a valid source overcover. For a real-valued
path $Y$ and integers $a\le z$, put
\[
 \Var(Y;[a,z])=\sum_{e=a}^{z-1}|Y(e+1)-Y(e)|,
 \qquad \Var(Y)=\Var(Y;[0,T-1]).
\]
In particular, $\Var(x)\le T-1$.

\begin{definition}[Complete families and the two-class condition]
\label{def:complete-families}
\formal{SourceCovers/FamilyAccounting.lean}{SourceCovers.actualOwner,SourceCovers.ActualFamily,SourceCovers.HasSpanColoring,SourceCovers.mem_family_edges_iff,SourceCovers.constant_extension_isWalk}
A long hop $(p,u)$ has $u-p\ge2$. Its side is upper or lower according
as the work head stays strictly above or strictly below its endpoint cell
between $p$ and~$u$. Its owner is the first time strictly between the
endpoints at which its distance from that cell is largest. A
\emph{complete actual family} consists of all long hops of one physical
head, on one side, with one owner. For a family $G$, let $o_G$ be that
owner, let $(p_G,U_G)$ be its outermost hop, let
$I_G=(p_G,U_G]$ be its full half-open span, and let $h_G$ be its hop count.
Families are formed from all actual long hops before any selection or
payment. They are never redefined by deleting inconvenient hops.

The full family system has an \emph{$N$-class span partition} if it
partitions into $N$ classes whose spans are pairwise disjoint within each
class. A selected core is \emph{Q2} if its complete spans admit such a
partition into two classes. Empty classes are permitted. The full-system
constant $N$ need not be two.
\end{definition}

These are the first-maximum families of
\cite[Sections~2.2 and~3.2]{EighthPaper2026}. In that paper a bounded
owner-load supplies a bounded number of disjoint-span classes. Here the
existence of a fixed $N$-class partition is an explicit hypothesis; we
do not infer it from the number of physical tapes alone.

\begin{definition}[Admissible return geometry]\label{def:admissible-geometry}
An admissible geometry with parameters $d,b,c,\beta,N,q,a_0$ consists
of a fixed complete geometry and nonempty correctness set $\mathcal C$
as above, with $T\le cL$, whole-step capacity at most $\beta$, an
$N$-class partition of its full actual family system, an original Q2
core $\mathcal K_0$, and a disjoint set $\mathcal E$ of at most $q$
other complete families. Every family outside
$\mathcal K_0\cup\mathcal E$ has at most $a_0$ hops.
The local computation is deterministic; no linearity or product
assumption is imposed on $\mathcal C$. The word ``admissible'' below
includes all these premises. In a sequence, $c,\beta,N,q$ are fixed;
the geometry, core and exceptional families may vary with dimension.
\end{definition}

\begin{lemma}[Family accounting]\label{lem:family-accounting}
\formal{SourceCovers/FamilyAccounting.lean}{SourceCovers.family_owner_inside,SourceCovers.family_hops_nested,SourceCovers.family_span_hop,SourceCovers.family_size_le_owner_arms,SourceCovers.two_mul_sum_family_size_le,SourceCovers.sum_family_variation_eq,SourceCovers.q2_accounting,SourceCovers.mem_boundaryLowers,SourceCovers.boundary_lowers_le_crossings,SourceCovers.boundary_lowers_le_family_sizes,SourceCovers.boundary_lowers_le}
The hops in a family are strictly nested. If its ends are
$u_1<\cdots<u_{h_G}=U_G$, their starts satisfy
$p_1>\cdots>p_{h_G}=p_G$, and every hop contains $o_G$ strictly.
The outer span is itself a hop, and
\[
 h_G\le\min(o_G-p_G,U_G-o_G)\le\frac{U_G-p_G}{2}.
\]
For any set $\mathcal K$ of complete families, let
$m_{\mathcal K}(e)=\card{\{G\in\mathcal K:p_G\le e<U_G\}}$
be its overlap at the elementary boundary $e\mid e+1$. For every real path $Y$,
\begin{align}\label{eq:family-overlap-profile}
 2\sum_{G\in\mathcal K}h_G
 &\le\sum_{e=0}^{T-2}m_{\mathcal K}(e),\notag\\
 \sum_{G\in\mathcal K}\Var(Y;[p_G,U_G])
 &=\sum_{e=0}^{T-2}m_{\mathcal K}(e)|Y(e+1)-Y(e)|.
\end{align}
In particular, a Q2 core has $m_{\mathcal K}(e)\le2$, and
\begin{equation}\label{eq:q2-mass-variation}
 \sum_{G\in\mathcal K}h_G\le T-1,
 \qquad
 \sum_{G\in\mathcal K}\Var(x;[p_G,U_G])
       \le2\Var(x)\le2(T-1).
\end{equation}
For a set $\mathcal U$ of unselected families, let $Z_{\mathcal U}(e)$
be the distinct lower endpoints of its work edges and the state edge
crossing $e\mid e+1$, and let
$c_G(e)=\card{\{(p,u)\in G:p\le e<u\}}$. At every boundary inside the
computation,
\begin{equation}\label{eq:local-boundary-bill}
 |Z_{\mathcal U}(e)|
 \le1+\sum_{\substack{G\in\mathcal U\\ p_G\le e<U_G}}c_G(e)
 \le1+\sum_{\substack{G\in\mathcal U\\ p_G\le e<U_G}}h_G.
\end{equation}
Thus, if every unselected family has at most $a$ hops, this count is at
most $1+m_{\mathcal U}(e)a\le1+Na$. The same conclusion holds under
the local bound $m_{\mathcal U}(e)\le N$, without supplying a span partition.
\end{lemma}

\begin{outline}
The formalization imports the eighth paper's actual-hop, owner and
family definitions and their verified nesting, outer-span and height
bounds. Uniqueness of previous visits makes nesting strict. After
reflection for a lower family, its hop count is at most the height
of the common owner above the outer endpoint cell. The unit-step
distance bound limits that height separately by $o_G-p_G$ and
$U_G-o_G$, giving the owner-sensitive estimate.

Index an elementary edge by its lower time $e$. The full span
$(p_G,U_G]$ contains its upper time $e+1$ exactly when
$p_G\le e<U_G$. Double counting these incidences gives
\cref{eq:family-overlap-profile}: use unit edge weights for span
lengths, and $|Y(e+1)-Y(e)|$ for variation. The shared counting
identity works for weights in any additive commutative monoid.
Families active at one boundary have distinct colors in any
disjoint-span coloring. A Q2 core therefore has overlap at most two;
the input's unit-step bound gives \cref{eq:q2-mass-variation}.

For the boundary count, the Lean set contains the state lower
endpoint and the union of all crossing-hop lowers in the active
families. Its membership theorem proves that this restriction to
active families omits no crossing hop. The finite-union and image
cardinality bounds give \cref{eq:local-boundary-bill}. Applying the
family-size bound $a$ and overlap bound $N$ gives $1+Na$.
This last counting argument needs neither unit steps nor nesting
once the local overlap and size bounds are supplied.
\end{outline}

The source-bank and grid-payment arguments below use the uniform
specializations. The exact overlap and crossing-hop profiles retain
information that can reduce a geometry's payment estimate; no stronger
uniform exclusion is inferred from them here.

\subsection{Clipped shadows and exact core orbits}

\begin{definition}[Paid clipping and selected orbits]\label{def:clipped-shadow}
\formal{SourceCovers/ClippedShadows.lean}{SourceCovers.ActualFamily.shadow,SourceCovers.ActualFamily.shadow_at_hop,SourceCovers.ActualFamily.shadow_exact,SourceCovers.shadowOrbit}
Pay the outer lower endpoint $p_G$ of each selected family. Its
\emph{clipped shadow} $s_G$ is the identity off $(o_G,U_G]$. On that
interval it initially holds $o_G$; at each hop end $u_i$ it takes the
value $p_i$ and holds it until the next end, with every held value
replaced by its maximum with $p_G+1$. For a family set $\mathcal S$,
let $\mathcal O_{\mathcal S}(t)$ be the set of values obtained from $t$
by all finite compositions of its clipped shadows, including the empty
composition.
\end{definition}

The shadow is antitone on its active interval, and every nonidentity
action satisfies
\begin{equation}\label{eq:strict-shadow}
 p_G<s_G(t)\le o_G<t\le U_G.
\end{equation}
Clipping changes an actual hop value only at the outermost end $U_G$,
whose lower endpoint has been paid. Hence every selected actual hop
with an unpaid lower endpoint is represented exactly. We do not gate
these maps at other paid vertices. Their orbits may retain extra paths
through paid vertices, which is safe for a source overcover and preserves
closure under each selected shadow.

For a path $Y$ in a metric space, use the same variation notation with
$\operatorname{dist}(Y(e),Y(e+1))$ in place of the absolute difference.
This agrees with the earlier definition for real paths.

\begin{lemma}[Variation of a shadow or successful path]
\label{lem:shadow-variation}
\formal{SourceCovers/ClippedShadows.lean}{SourceCovers.jointShadow_eq_on,SourceCovers.jointShadow_variation,SourceCovers.shadow_success_cons_iff,SourceCovers.shadow_word_structure,SourceCovers.shadow_word_variation,SourceCovers.variation_comp_le_crossings,SourceCovers.exists_sharp_patch}
The joint map of clipped shadows with pairwise disjoint full spans has
\[
 \Var(Y\circ s)\le3\Var(Y)
\]
for every metric-valued path $Y$. For any fixed word of clipped shadows,
the set of starting times at which every action is nonidentity is an
integer interval, possibly empty. Its endpoint map is monotone in one
direction on that interval and no greater than its argument. Extending
that endpoint map by the identity outside the success interval also
gives variation at most $3\Var(Y)$ for every $Y$. The word may be empty
and may repeat families.

More generally, for a time map $f:\{0,\ldots,n\}\to\{0,\ldots,N\}$,
let $c_f(z)$ count the indices $e<n$ for which
$\min(f(e),f(e+1))\le z<\max(f(e),f(e+1))$. Then
\begin{equation}\label{eq:crossing-variation}
 \Var(Y\circ f;[0,n])
 \le\sum_{z=0}^{N-1}c_f(z)\operatorname{dist}(Y(z),Y(z+1)).
\end{equation}
The factor three is attained in the general class of regressive interval
patches, even by one constant active map.
\end{lemma}

\begin{outline}
The concrete held value is the minimum, over all hops in a complete
family, of the hop's start if it has ended and the owner otherwise.
It is antitone in time. Strict nesting identifies this minimum with
the actual start at each hop end. Taking its maximum with $p_G+1$
gives exactly the clipped shadow: every hop whose lower endpoint is
unpaid is preserved. Its active interval is exactly its nonidentity set.

For the reusable variation argument, sum the triangle inequality along
the elementary edges crossed by each jump, then interchange the finite
sums. The shared incidence identity gives \cref{eq:crossing-variation}.
A map monotone in either direction on one interval and equal to the
identity elsewhere has at most one downward crossing of each threshold.
Indeed, two downward crossings would force the intervening points into
the active interval. An antitone map could not rise between them; for a
monotone map, both downward crossings would have to occur at the two
joins, contradicting the order of the exterior identity values.
Telescoping the threshold's zero--one indicator bounds the total number
of crossings by twice the downward count plus one, hence by three.
This argument needs only preservation of the finite time range, not
regression of the patch.

For a disjoint packing, a threshold outside a family's full span is
unaffected: both the argument and its image stay on the same side.
At most one full span straddles a given threshold, so the packed map
has the same threshold indicator as that one shadow, or as the identity.
The three-crossing bound therefore holds independently of the family count.

Induct on a word. Its success domain is the first action's interval
intersected with the preimage of the tail's success domain. Mathlib's
order-convex preimage lemmas make this an interval. Composing the
monotone restrictions preserves one of the two orientations, and every
action weakly decreases time even when inactive. The successful endpoint
map therefore has the required interval, monotonicity and range properties.
Apply the single-patch bound to its identity extension. Empty success
domains and the empty word are included.

For sharpness of the general patch bound, take $f(t)=1$ on $(2,4]$ and
$f(t)=t$ elsewhere, on times $0,\ldots,5$. With $Y(t)=0$ for $t<2$ and
$Y(t)=1$ otherwise, $\Var(Y)=1$ and $\Var(Y\circ f)=3$. Lean
checks this finite witness directly. No sharpness claim for a particular
machine's source bank is inferred from this example.
\end{outline}

This lemma supplies the any-argument-path bound used repeatedly in the
exact source bank, without a factor exponential in successful word length.
The crossing profile can improve a particular map's estimate; the later
uniform constants still use the factor-three specialization. The packed
bank is formalized next; recovery through the residual dependencies remains
the subsequent task.

\begin{lemma}[The Q2 overlap forest and protected phases]
\label{lem:protected-phase}
\formal{SourceCovers/ProtectedPhases.lean}{SourceCovers.q2_caterpillar,SourceCovers.protected_survives,SourceCovers.admissiblePhases_card,SourceCovers.phase_payment_sum,SourceCovers.phase_average_exact,SourceCovers.protected_phase_average,SourceCovers.nonleaf_mass_le_total,SourceCovers.phase_path_length,SourceCovers.phase_success_length,SourceCovers.phase_shadowOrbit_exact,SourceCovers.phase_orbit_chains,SourceCovers.phase_orbit_card}
Join core families whose half-open spans intersect. For a Q2 core this
graph is a caterpillar forest: the vertices with at least two neighbors,
called the \emph{original nonleaves}, form a path within each component
where they occur. Let $\mathcal N$ be their set. Rank these paths from
left to right, starting at one separately in each component. All these
choices refer to the core before deletion.

Fix $M\ge2$ and a protected set $\mathcal P$ of at most $R<M$ core
families. Let $\mathcal A$ be the residues modulo $M$ not used by the
ranks of $\mathcal P\cap\mathcal N$, and let $F$ be the number of
distinct residues they use. For $\rho\in\mathcal A$, delete only the
original nonleaves of rank congruent to $\rho$, paying every actual
lower endpoint of their hops; call this union $\Gamma_{\rm rank}(\rho)$.
Every protected family survives. Put
$\tau_{\mathcal N}=\sum_{G\in\mathcal N\setminus\mathcal P}h_G$.
Then
\begin{equation}\label{eq:rank-phase-average}
 \card{\mathcal A}=M-F\ge M-R,
 \qquad
 \frac1{\card{\mathcal A}}
 \sum_{\rho\in\mathcal A}\card{\Gamma_{\rm rank}(\rho)}
 \le\frac{\tau_{\mathcal N}}{\card{\mathcal A}}
 \le\frac{\tau_{\mathcal N}}{M-R}.
\end{equation}
More precisely, the sum of paid-union sizes is at most
$\sum_{G\in\mathcal N\setminus\mathcal P}\card{L_G}$, where
$L_G$ is the set of actual hop lowers of $G$. The average bound with
denominator $M-F$ only requires $F<M$, regardless of $\card{\mathcal P}$.

In every phase, each surviving simple overlap path has at most $M+1$
vertices; in particular each successful core word has at most $M+1$
actions. The complete core orbit is the union of at most two deterministic
chains from the root and has at most $\ell=2M+3$ time values. Original
leaves and components of one or two vertices need not be deleted.
\end{lemma}

\begin{outline}
The overlap graph has no triangle: three pairwise intersecting intervals
have a common point, but one Q2 class would contain two of them. It has
no longer cycle either. On a cycle choose a vertex with largest left
endpoint. Its two cycle neighbors both extend strictly past that
endpoint, hence overlap each other, giving a triangle. Thus the graph
is a forest. Lean reuses Mathlib's cycle rotation and path operations.

An interval contained in another can have no further neighbor without
creating a triangle. Consequently nonleaf intervals have distinct left
endpoints and strictly increasing right endpoints in that order. An
overlap edge between nonleaves cannot skip a third nonleaf in the order:
that third interval would meet both endpoints and create a triangle.
Lean defines the rank as one plus the number of earlier nonleaves in
the same connected component, and proves that adjacent ranks differ
by one and distinct nonleaves have distinct ranks. Every internal vertex
of a simple path is a nonleaf, so removing original leaves preserves
connectivity between nonleaves. These facts give the connected nonleaf
path and its maximum degree two.

A finite unit-step integer chain avoiding one residue modulo $M$ visits
at most $M-1$ distinct integers. Lean proves this using the first forbidden
integer at or above the chain's minimum; a unit step cannot cross that
integer without visiting it. Apply this to the internal vertices of a
surviving simple path. Their ranks are distinct, so including the two
endpoints gives at most $M+1$ vertices. This argument needs no deletion
of leaves and also covers components of one or two vertices.

By \eqref{eq:strict-shadow}, consecutive successful actions have
overlapping spans and strictly decreasing owners. Their names form a
simple path, to which the preceding bound applies.
After an action from one Q2 class, its image lies in the same family's
span at or before its owner. That family is now inactive and no other
family of its class can act there. The only possible next action is
therefore the unique active family of the other class. At the root
there are at most two choices of first action; thereafter continuation
is deterministic. Including all intermediate stopping points gives
at most $1+2(M+1)=2M+3$ values.

The shared finite-orbit proof erases inactive actions from an arbitrary
word without changing its endpoint. It identifies the full word orbit
exactly with the values reachable in at most $M+1$ nonidentity steps.
This includes the empty word, intermediate stopping points and any
collisions between the two chains.

Only protected nonleaves forbid residues. Each unprotected nonleaf is
deleted in at most one admissible phase. Mathlib's finite-union bound
charges its actual lower support, whose size is at most its hop count.
Interchanging the finite sums gives the payment inequality, and counting
distinct forbidden residues gives $\card{\mathcal A}=M-F$.
\end{outline}

This supplies the finite core-orbit bound needed for the exact source
bank. Compared with deleting leaves together with their neighbors, it
retains the same orbit bound while paying only for original nonleaves.
The earlier estimate using
$\tau=\sum_{G\in\mathcal K\setminus\mathcal P}h_G$ still follows from
$\tau_{\mathcal N}\le\tau$. The next theorem assembles these bounds into a uniform packed bank.

\begin{theorem}[Theorem S: an exact source bank]\label{thm:exact-source-bank}
\formal{SourceCovers/ExactSourceBank.lean}{SourceCovers.exact_source_bank,SourceCovers.exactSourceBank_length,SourceCovers.exactSourceBank_bounded,SourceCovers.exactSourceBank_exact,SourceCovers.exactSourceBank_cover,SourceCovers.exactSourceBank_raw_root,SourceCovers.exactSourceBank_raw_values_card,SourceCovers.exactSourceBank_raw_selected_edge,SourceCovers.exactSourceBank_raw_variation}
Select a Q2 core of at most $J$ complete families, at most $q$ further
complete exceptional families, and a protected subcore of size at most
$R$. Let $M\ge\max(2,R+1)$ and use any admissible phase of
\cref{lem:protected-phase}. Pay the rank lowers and the outer lower
endpoints of all selected core and exceptional families. Put
\begin{equation}\label{eq:bank-constants}
 \begin{split}
  \ell&=2M+3,\qquad m_0=4+2M(M-1),\qquad C_0=16+24M(M-1),\\
  m_q&=m_0(1+m_0)^q,\qquad C_q=C_0(1+3C_0)^q.
 \end{split}
\end{equation}
There is a list $\mathcal B$ of at most $m_q$ geometry-fixed total time
maps such that, simultaneously for all time vertices $t$,
\begin{enumerate}
\item $\{a(t):a\in\mathcal B\}$ is exactly the complete orbit of the
      retained core and all selected exceptional shadows;
\item for every real-valued path $Y$,
      $\sum_{a\in\mathcal B}\Var(Y\circ a)\le C_q\Var(Y)$;
\item the time-value set is a union of at most $(\ell+1)^q$ complete
      core orbits, each a union of the two chains in
      \cref{lem:protected-phase}.
\end{enumerate}
Consequently the raw bank $Y_a=x\circ a$ contains the root read,
has total variation at most $C_q(T-1)$, and is set-closed across every
selected actual edge whose lower endpoint is unpaid. Its pointwise
cardinality is at most $\ell(\ell+1)^q$. For fixed $M$ and $q$, the
constants do not depend on $J$ or on the sizes or scales of the original
families.
\end{theorem}

\begin{outline}
Keep the original nonleaves as spine vertices. In each two-vertex
component, add its color-zero vertex to the spine; leave isolated vertices
in the leaf pack. The remaining leaves have pairwise disjoint spans,
so their joint shadow $s$ has variation constant three.
Every internal vertex of a successful word is an original nonleaf.
Thus packing leaves occur only at the two ends of such a word.

Label a surviving nonleaf by its original rank modulo $M$, excluding
the deleted residue. These $H=M-1$ labels are injective within each
surviving spine component: along a path, ranks make unit steps and avoid
the deleted residue, so they stay in an interval shorter than one period.
Equal residues therefore mean equal ranks and equal vertices. The extra
spine vertices from two-vertex components are isolated in the spine and
may use any of the $H$ labels. This uses $M\ge2$.

For each unordered pair of labels, with repetition allowed, pack one map
across all surviving spine components. In a component, choose the
nonempty successful word with those endpoint labels, if one exists.
The forest has a unique simple path between two vertices, and strictly
decreasing owners fix its orientation. Hence that word is unique.
Use its endpoint map on its success interval and identity elsewhere.
If no word exists, use identity throughout that component.

The union of the spans in a connected component is its interval hull,
and distinct component hulls are disjoint. Each chosen word remains in
its component hull. The shared disjoint-patch theorem therefore bounds
each packed map's variation by three, including both joins to identity.
Add an identity slot. Erasing inactive actions from any spine word shows
that these slots give exactly its complete orbit, simultaneously at every
root. Their number and summed variation coefficient are
\[
 m_S=1+\binom M2,
 \qquad C_S=1+3\binom M2.
\]
All choices precede the root and the argument path; a slot's representing
word is allowed to vary between components.

For each spine slot $a$, use
\[
 a,\qquad a\circ s,\qquad s\circ a,\qquad s\circ a\circ s.
\]
The two endpoint leaves and the spine middle of every successful core
word are represented by these four choices. Conversely each value follows
a core word at that argument. This proves exactness, rather than merely
coverage. The any-path variation bound composes: the four contributions
are at most $C_S,3C_S,3C_S,9C_S$. Thus there are $4m_S=m_0$
slots with total coefficient $16C_S=C_0$. The pointwise bound remains
$\ell$, by \cref{lem:protected-phase}.

Sort the exceptional families by decreasing owner, retaining a fixed tie
order. Successful words have strictly decreasing owners, so each such word
uses an ordered exceptional subsequence, with core words between its
members. Let $B$ be the core bank and $D$ the bank for the remaining
exception list. On inserting its first shadow $f$, form
\[
 D\ \mathbin{+\!+}\ \{b\circ f\circ a:b\in D,\ a\in B\},
\]
where $+\!+$ denotes list concatenation. Exactness follows by splitting a
successful word at $f$, or omitting $f$ if it is absent. Every constructed
value follows a selected word, even though a packed core slot need not
represent one fixed word globally. Slot count and summed variation obey
\[
 m_{r+1}=(1+m_0)m_r,
 \qquad C_{r+1}=(1+3C_0)C_r.
\]
These recurrences give \eqref{eq:bank-constants} without separately
formalizing the binomial sums over subsequences.

For the complete-core-orbit cover, retain the starting vertices of each
final core continuation. Each inserted exception creates at most $\ell$
new choices for each old continuation pattern, in addition to the pattern
omitting it. Induction gives at most $(\ell+1)^q$ terminal core roots.
The equality of the covered set with the enlarged orbit proves its
pointwise cardinality bound.

The empty word gives the root read. Prefixing a selected shadow gives
orbit set-closure. For an original core family removed by the phase,
every actual lower endpoint lies in the rank payment; it therefore has
no unpaid actual edge. For every remaining selected family, strict
clipping represents the actual hop exactly unless its outer lower was
paid. This proves closure across all originally selected unpaid edges,
and taking images under $x$ preserves it. Finally, the unit-step bound
$\Var(x)\le T-1$ gives the raw variation estimate. Other paid vertices
do not truncate the orbit.
\end{outline}

The formal variation bound holds for every pseudometric-valued path and
every finite horizon, with the same constants. The geometric bank needs
no unit-step assumption on the work walks; that assumption is used when
identifying shadows with actual hops. It also needs no numerical bound
on the number of core families. In the generic exceptional-list argument,
repeated entries and overlap with the core are allowed; the displayed
bounds count list entries.


Theorem~S supplies a representation of complete selected dependence
whose variation cost, for fixed $M$ and $q$, is independent of the number
of selected core families.
The next lemma pays for the remaining actual dependencies. Without this
step, Theorem~S alone is not a transposition certificate.

\subsection{Residual cuts and complete functional recovery}

For an integer width $w\ge1$ and shift
$\theta\in\{0,\ldots,w-1\}$, define the linear, endpoint-clipped block
\begin{equation}\label{eq:raw-block}
 B_\theta(y)=
 \{z\in\Omega:\lfloor(z+\theta)/w\rfloor
                   =\lfloor(y+\theta)/w\rfloor\}.
\end{equation}
Thus $|z-y|\le w-1$ within a block, including a partial block at either
end of the input.

\begin{lemma}[Paid residual-source recovery]\label{lem:residual-source-recovery}
\notformalized
Let $\Gamma_0$ be an initial paid set, and let $Y_1,\ldots,Y_m$ be
geometry-fixed maps from the timeline to $\Omega$. Their value set
$A(t)$ contains $x(t)$. Partition all dependency edges into selected
and residual edges. Suppose $A(u)\subseteq A(v)$ for every selected
edge $v\to u$ with $u\notin\Gamma_0$. For an unpaid lower endpoint
$u$, let $m_R(u)$ be its latest residual successor, or $u$ if none.
Define
\[
 V=\sum_i\Var(Y_i),\qquad
 \mathcal R=\sum_{u\notin\Gamma_0}\sum_i
   \operatorname{osc}(Y_i;[u,m_R(u)]),
\]
where oscillation is maximum minus minimum. If at most $H$ distinct
unpaid residual lowers are live across every elementary time boundary,
then, for every $w\ge1$, there are actual residual-lower sets
$\Gamma_\theta$ satisfying
\begin{equation}\label{eq:residual-average}
 \mathbb E_\theta\card{\Gamma_\theta}
       \le\frac{\mathcal R}{w}\le\frac{HV}{w}.
\end{equation}
For each shift, after paying
$\Gamma=\Gamma_0\cup\Gamma_\theta$, every output at time $t_r$ is
determined on $\mathcal C$ by the same transcript $G_\Gamma$ and the
bits in
\begin{equation}\label{eq:complete-bank-recovery}
 \bigcup_{i=1}^m B_\theta(Y_i(t_r)).
\end{equation}
No bound on full path depth, linearity or independence on a genie fibre
is needed.
\end{lemma}

\begin{informalproof}
For each unpaid residual lower $u$, pay $u$ if, for some $i$, the label
$\lfloor(Y_i(t)+\theta)/w\rfloor$ is not constant on
$[u,m_R(u)]$. This is an actual union of lower endpoints, so a lower
used by several dependencies is paid once. If the integer range of one
slot on an interval has oscillation $z$, a uniformly shifted grid cuts
that range with probability $\min(1,z/w)$. The union bound over slots
and lowers gives the first inequality of \eqref{eq:residual-average}.
Oscillation is at most variation on the interval. Summing that variation
over lower lifetimes charges an elementary step at most $H$ times,
giving $\mathcal R\le HV$.

A surviving residual edge preserves every slot label. A surviving
selected edge preserves inclusion of the sets of labels by the
hypothesis. Along any surviving backward path, induct from its last
vertex to its root: the label of its direct input read belongs to the
root's label set, because $x(t)\in A(t)$ at the last vertex. Therefore
every full surviving raw source belongs to
\eqref{eq:complete-bank-recovery}. This treats paths of arbitrary
length and all alternations of state and physical work edges.

Finally recursively evaluate the deterministic step functions. A
parent in $\Gamma$ is supplied by $G_\Gamma$; every other parent is
evaluated from its own parents and direct read. Time strictly decreases,
so this recursion terminates. Initial values and reads outside the input
are constants. All unfixed raw reads belong to the displayed union,
which therefore determines the output. The same reasoning starts at
a revealed output root and retains its direct read and all unpaid
parents; revealing the root was never used to delete them.
\end{informalproof}

For Theorem~S, take the surviving selected family edges as selected,
retain all other state and physical work edges as residual, and disregard
only edges already blocked at their lower endpoints. If every family
outside the selected core and the selected exceptional set has at most
$a$ hops,
\cref{lem:family-accounting} gives $H\le1+Na$. In each admissible phase
there is a shift with
\begin{equation}\label{eq:native-genie-ledger}
 \card\Gamma\le
 J+q+\card{\Gamma_{\rm rank}(\rho)}
       +\frac{(1+Na)C_q(T-1)}w.
\end{equation}
In particular a phase attaining the average in
\eqref{eq:rank-phase-average} permits replacing its rank payment by
$\tau_{\mathcal N}/(M-R)$, or the older upper bound $\tau/(M-R)$.
More generally, one may average a complete signed
reward-minus-payment expression over the admissible phases and their
raw-grid shifts. A reward chosen in one phase or shift cannot be paired
with a payment from another. Deleted selected families are blocked by
their paid lowers; they are not silently reclassified as small residual
families.

\subsection{Cyclic freezing of actual transpose demands}

View $\Omega$ as the integer points on a circle of circumference $L$.
For a nonempty subset $S$, list the circular distances between consecutive
distinct points as $g_1,\ldots,g_v$; their sum is $L$. A singleton has
one gap of length $L$. Define its point-gap area by
\begin{equation}\label{eq:gap-area}
 \gaparea(S)=\frac1{2L^2}\sum_{h=1}^v g_h^2.
\end{equation}
The blocks in \eqref{eq:raw-block} remain linear blocks. The circle is
used only for an overcover and a common ordering of raw coordinates.

\begin{theorem}[Theorem F: cyclic freezing]\label{thm:cyclic-freezing}
\notformalized
Let $m\ge1$ and let $Y_i:\{0,\ldots,T-1\}\to\Omega$, $1\le i\le m$,
be geometry-fixed bank maps. Let $Z$ be a
deterministic helper determined by the input, with $H(Z)\le g$ for the
uniform law on $\mathcal C$. Suppose every output is determined on
$\mathcal C$ by $Z$ and the raw block union in
\eqref{eq:complete-bank-recovery}, at one width $w$ and shift $\theta$. Put
$V=\sum_i\Var(Y_i)$. For output column $j$, let
$s_j=t_{jW}$ be its first append time and set
$S_j=\{Y_i(s_j):1\le i\le m\}$ and $A_j=\gaparea(S_j)$. Then
\begin{equation}\label{eq:cyclic-freezing}
 \Delta\ge
 W\sum_{j=0}^{d-1}A_j
       -\frac{2V}{d}-2(w-1)-\frac{W+1}{2}
       -g.
\end{equation}
In particular the bound holds with $g=\beta\card\Gamma$ for an actual
whole-step genie. If $u$ additional raw bits at geometry-fixed coordinates
are revealed jointly with that genie, it holds with
$g=\beta\card\Gamma+u$.
The bound permits arbitrary nonlinear recovery and a correlated
correctness set. Its area term uses the actual transpose demand stripes,
not an assumption that demands and sources are independent.
\end{theorem}

\begin{informalproof}
Write $e_j=t_{(j+1)W-1}$ and
$\delta_j=\sum_i\Var(Y_i;[s_j,e_j])$. These column intervals have
disjoint elementary edges, so $\sum_j\delta_j\le V$. Every source
block used during column $j$ lies in the fixed circular overcover
\[
 T_j=\{z\in\Omega:\dist_{\rm circ}(z,S_j)\le r_j\},
 \qquad r_j=w-1+\delta_j.
\]
Indeed, its coordinate differs from the corresponding current anchor
by at most $w-1$, and that anchor differs from its frozen value by at
most $\delta_j$. This covers entire blocks, not just their representatives.

Choose a uniformly random circular cut and an independent uniformly
random orientation. They define one total order on all raw coordinates,
shared by every column. Select a demanded coordinate $z$ of column $j$
when $z\notin T_j$ and $z$ precedes every coordinate in $T_j$. Every
recovery edge between selected coordinates then points forward in this
same order, so their induced recovery relation is acyclic.

Consider a complementary circular run of $T_j$ having $a$ integer
points. For its $v$-th point, the two orientations give selection
probabilities $v/L$ and $(a+1-v)/L$. Their average is the constant
$(a+1)/(2L)$ throughout the run. The selection probability is zero
on $T_j$. The sum of these constants over complementary runs is at
most $1/2$: the run points together with one distinct covered endpoint
per run number at most $L$. If $T_j=\Omega$, all these sums are empty.

For a frozen point gap of length $g$, thickening by integer radius $r_j$
leaves $(g-1-2r_j)_+$ integer points. The mean selection probability
over all $L$ raw coordinates is consequently
\[
 A_{j,r_j}=\frac1{2L^2}
       \sum_{g}(g-1-2r_j)_+\bigl((g-1-2r_j)_++1\bigr).
\]
For integers $g\ge1,r\ge0$,
\[
 g^2-(g-1-2r)_+\bigl((g-1-2r)_++1\bigr)\le(4r+1)g.
\]
When $g\ge2r+1$, expansion proves the inequality; otherwise its
left side is $g^2\le(4r+1)g$. Summing and using $\sum g=L$ gives
\begin{equation}\label{eq:thickening-area}
 A_{j,r_j}\ge A_j-\frac{2r_j+1/2}{L}.
\end{equation}
There is no factor for the number of slots or distinct frozen points.

The actual demanded coordinates in column $j$ form the stripe
\[
 D_j=\{iW+jb+k:0\le i<d,\ 0\le k<b\}.
\]
For every circular integer interval $I$,
\begin{equation}\label{eq:stripe-discrepancy}
 \left|\card(D_j\cap I)-\frac{\card I}{d}\right|
       \le b(1-1/d)\le b.
\end{equation}
To see this, the indicator of $D_j$ has period $W$ and one occupied
block of length $b$ per period. Its centered prefix count
$F(n)=\card(D_j\cap[0,n))-n/d$ is periodic and has oscillation
$b(1-1/d)$: it rises with slope $1-1/d$ on the occupied block and
falls with slope $-1/d$ on its complement. An interval discrepancy
is a difference of two prefix values, including intervals across the
circular seam.

Apply \eqref{eq:stripe-discrepancy} to the complementary runs, weighted
by their constant selection probabilities. Their constants sum to at
most $1/2$, so the expected number of selected demands in column $j$
is at least
\[
 W A_{j,r_j}-b/2
 \ge W A_j-\frac{2r_j}{d}-\frac1{2d}-\frac b2.
\]
Summing over columns and using
$\sum_j r_j\le d(w-1)+V$ gives expected selected size at least
the first four terms on the right of \eqref{eq:cyclic-freezing}.

For any realized common order, let $A$ be its selected coordinates.
Given $Z$ and the raw complement $\Omega\setminus A$, recover
the selected coordinates in reverse order using the acyclic relation.
For $X$ uniform on $\mathcal C$, it follows that
\[
 \log_2\card{\mathcal C}=H(X)
       \le H(Z)+H(X_{\Omega\setminus A}\mid Z)
       \le g+L-\card A.
\]
Average this inequality over the common cut and orientation to obtain
\eqref{eq:cyclic-freezing}. This entropy argument does not infer a
transcript cardinality bound from an entropy bound, and requires no
product distribution on a helper fibre. The genie and joint-revelation
special cases follow from subadditivity of entropy and the stated bit
capacities.
\end{informalproof}

Theorem~F charges the same genie that supplies the complete bank.
Own-source demands are rejected in its acyclic selection, so no separate
own-block collision loss is added. If $M,q,c$ are fixed, Theorem~S gives
$V=O(L)$, and the variation and stripe errors are $o(L)$ as
$d\to\infty$. The substantive remaining question is whether the
point-gap area exceeds the rank and residual payments; the following
sections prove this under explicit source-cover hypotheses.

\section{Fractional source covers and a paid deficit}\label{sec:source-covers}

We now improve the geometric term in \cref{thm:cyclic-freezing}.
The recovery bank remains the complete bank of \cref{thm:exact-source-bank}.
The covers introduced here are witnesses used to estimate its raw source
sets; they are not additional records available to the machine.

\subsection{Parameters and complete core covers}

Fix an admissible geometry, a selected $\mathrm Q_2$ subcore of at most $J$
families, and at most $q$ complete exceptional families. Choose integers
$K,R\ge0$, $J\ge\max(K,R)$ and $M\ge\max(2,R+1)$, distinguish at most $K$ selected
core families, protect at most $R$ selected core families, and choose a legal rank phase
$\rho$. All distinguished families are among those selected before the
phase; they need not be protected. Write
\begin{equation}\label{eq:cover-parameters}
 \ell=2M+3,\qquad C_0=16+24M(M-1),\qquad C_q=C_0(1+3C_0)^q.
\end{equation}
Let $A_\rho(u)$ be the complete orbit of the retained core at time $u$,
and $O_\rho(t)$ the complete orbit with the selected exceptions included.
The raw source set is $S_\rho(t)=x(O_\rho(t))$. Put
\begin{equation}\label{eq:alpha-tau}
 \alpha L=\sum_{G\text{ selected core, not distinguished}}
              \Var(x;[p_G,U_G]),\qquad
 \tau=\sum_{G\text{ selected core, not protected}} |G|.
\end{equation}
These sums use original full spans and hop counts, before rank deletion;
in particular they do not depend on $\rho$.

\begin{definition}[Fractional complete-core cover]\label{def:fractional-core-cover}
\formal{SourceCovers/CoreFractionalCover.lean}{SourceCovers.coreRawSets,SourceCovers.fractionalCoreCover,SourceCovers.fractionalCoreCover_exists_optimal,SourceCovers.one_le_fractionalCoreCover_toReal,SourceCovers.coreRawSets_phase_shadowOrbit_eq,SourceCovers.fractionalCoreCover_ne_top_iff_subset,SourceCovers.fractionalCoreCover_exactSourceBank_ne_top}
For a nonempty finite set $S_0\subseteq\{0,\ldots,L-1\}$, define
$\kappa_\rho(S_0)$ as the minimum of
\[
 \sum_{u=0}^{T-1}\lambda_u
 \quad\text{subject to}\quad
 \lambda_u\ge0,\qquad
 \sum_{u:\,a\in x(A_\rho(u))}\lambda_u\ge1\quad(a\in S_0).
\]
If no cover exists, set $\kappa_\rho(S_0)=+\infty$.
Covering orbits may be rooted at any computation time; their raw sets
may contain points outside $S_0$. A finite cover has an optimum, since
weights may be truncated at one and the resulting feasible set is
compact. It has weight at least one. In fact a finite cover exists
exactly when $S_0\subseteq\{x(u):0\le u<T\}$: every complete core
orbit contains its root and only earlier times. Thus unit weights at
all computation times suffice whenever the condition holds, including
for $S_0=S_\rho(t)$ at every computation time $t$.
\end{definition}

\begin{lemma}[Fractional interval covers]\label{lem:interval-rounding}
\formal{SourceCovers/IntervalCover.lean}{SourceCovers.exists_cover_packing,SourceCovers.exists_subcover_le_weight}
Let $S$ be a finite subset of a linearly ordered set, and let
$(A_i)_{i\in I}$ be a finite family of order-convex sets covering $S$.
Order-convex means that each set contains every point between two of
its points. There exist $C\subseteq I$ and $P\subseteq S$ such that
\[
 S\subseteq\bigcup_{i\in C}A_i,\qquad |C|=|P|,\qquad
 |P\cap A_i|\le1\quad(i\in I).
\]
Consequently, nonnegative weights $w_i$ satisfying
$\sum_{i:\,x\in A_i}w_i\ge1$ for every $x\in S$ admit an integral
subcover of cardinality at most $\sum_{i\in I}w_i$.
The target $S$ may be empty. In particular, this applies to closed,
open or half-open real intervals without a boundedness assumption.
\end{lemma}
\begin{outline}
Induct on the finite target. For a nonempty target let $x$ be its
least point, and let $y$ be the greatest target point sharing an
$A_i$ with $x$. Choose such a set $A_j$ containing $x,y$. Order
convexity makes $A_j$ cover all target points through $y$. No set
containing $x$ meets the residual target above $y$, by the choice of
$y$. Apply induction there after removing $j$ from the index family,
then add $j$ to the subcover and $x$ to the packing. Both cardinalities
increase by one, and every original set still contains at most one
packing point. The empty target uses two empty sets.

Fractional coverage implies ordinary coverage. Sum its inequalities
over the packing and interchange the finite sums. Each set's weight
occurs at most once, so $|C|=|P|\le\sum_iw_i$. The proof uses only
finite extrema, order convexity and finite sums; the shared library
supplies the declarations, which this paper re-exports.
\end{outline}
This rounding step preserves the total fractional weight in the
raw-source gap bound below. It does not alter the complete recovery bank.

\begin{lemma}[A fractional raw-source gap bound]\label{lem:fractional-gap}
\formal{SourceCovers/FractionalGap.lean}{SourceCovers.fractional_phase_gap,SourceCovers.fractional_core_gap,SourceCovers.fractional_core_gap_of_union_variation,SourceCovers.fractional_core_gap_of_weights}
With the parameters above, any nonempty $S_0$ with finite
$\kappa_\rho(S_0)$ satisfies
\begin{equation}\label{eq:fractional-gap}
 \gaparea(S_0)\ge
 \frac{(1-\alpha)_+^2}{2(2K+1)\kappa_\rho(S_0)},
 \qquad (z)_+=\max(z,0).
\end{equation}
Here $\gaparea(S)=\sum_g |g|^2/(2L^2)$, with the sum over the circular
gaps of the distinct points in $S$; a singleton has one gap of length $L$.

A sharper form retains actual component counts. Join consecutive
computation times across the elementary edges in the union $B$ of the
original nondistinguished full spans. Let $m_\rho(u)$ count the resulting
time components met by $A_\rho(u)$ and put
$V_B=\sum_{e\in B}|x(e+1)-x(e)|$. For any feasible cover weights,
\begin{equation}\label{eq:fractional-gap-incidence}
 \gaparea(S_0)\ge
 \frac{(1-V_B/L)_+^2}{2\sum_{u=0}^{T-1}m_\rho(u)\lambda_u}.
\end{equation}
Here $V_B\le\alpha L$ and $m_\rho(u)\le2K+1$.
The same conclusions hold for any retained subset of the selected
$\mathrm Q_2$ core. No rank phase or unit-speed assumption on $x$ is
needed for this estimate.
\end{lemma}
\begin{outline}
Label a time $t$ by the number of elementary edges before $t$ that are
outside $B$. Equal labels define contiguous time components, including
isolated times. Their internal edge sets partition $B$ within the
computation horizon. The interval
between the minimum and maximum raw coordinates on each component has
length at most its internal variation. Thus the sum of all component
interval lengths is at most $V_B$, with every elementary edge charged once.

A nondistinguished shadow stays in one component, since every edge
between its endpoints belongs to its original full span. A complete
core orbit has at most two initial successors and thereafter deterministic
continuation. Along either branch, a label change uses a distinguished
family whose owner precedes the current time. The new time is at most
that owner, so this family cannot pay for another change on the branch.
Counting the shared root label once gives $m_\rho(u)\le1+2K$.
The formal proof keeps the smaller count of retained distinguished
owners preceding the root before relaxing it to $K$.

Give each component interval the sum of the weights of the covering
orbits that meet it. These weights cover $S_0$ and have total
$\sum_um_\rho(u)\lambda_u$. By \cref{lem:interval-rounding}, they
admit an integral subcover with at most that many intervals. Their
total length is at most $V_B$.

Sort the distinct target points and append the first point plus $L$.
Call a consecutive gap exposed if no chosen interval contains both
endpoints. One interval can contain the left endpoint of at most one
exposed gap: if it contained two, it would contain both endpoints of
the earlier gap. Since every left endpoint is covered, the exposed
gap count is at most the subcover size. The other gaps have disjoint
elementary edge sets inside the covering intervals, so their total
length is at most $V_B$. The exposed gaps therefore have total length
at least $(L-V_B)_+$. Cauchy--Schwarz, followed by division by $2L^2$,
gives \cref{eq:fractional-gap-incidence}. The unrolled last gap handles
the circle seam and singleton target without an extra component.

Finally take an attained minimum cover, use $V_B\le\alpha L$ and
$m_\rho(u)\le2K+1$, and obtain \cref{eq:fractional-gap}. All extrema,
gap lengths and component incidences in this argument are constructed
from the actual time path and the complete core orbits.
\end{outline}
This supplies a verified geometric credit for the paid-deficit argument.
The shared-root count improves the former factor $4(K+1)$ to $2(2K+1)$;
the quantitative consequences below use the improved factor. It does not
supply economical covers for arbitrary native geometries.

\subsection{Anchored approximation}

We represent finite circular geometry by a common subdivision containing
all points and arc endpoints. Its elementary edges have real lengths
$w_e\ge0$ summing to $L>0$. A closed arc contains consecutive edges and
both endpoint vertices; its union length is the sum of the lengths of
the edges in that union. Thus an edge covered several times is counted
once. Subdivision indices are labels, not physical integer coordinates.
The formalization identifies each weighted gap with its real endpoint
displacement, including the gap across the chosen seam. Taking all edge
lengths equal to one recovers the raw-coordinate definition used above.
Write $\psi(e)=e-e^2/2$.

\begin{lemma}[Anchored erosion of circular gaps]\label{lem:anchored-erosion}
\formal{SourceCovers/AnchoredErosion.lean}{SourceCovers.penalty, SourceCovers.area_lower_bound, SourceCovers.area_lower_bound_linear, SourceCovers.anchored_displacement_area_lower_bound, SourceCovers.anchored_point_area_lower_bound}
Let $S_0,S$ be nonempty finite point sets on the circle of circumference
$L$, represented in a common subdivision as above. Suppose that
$S\subseteq S_0\cup\bigcup_{i=1}^m I_i$, where the $I_i$ are closed
circular arcs, each containing a point of $S_0$. If their union has
length at most $E$, set $e=\min\{E/L,1\}$. Then
\begin{equation}\label{eq:anchored-erosion}
 \gaparea(S)\ge\gaparea(S_0)-\psi(e)
             \ge\gaparea(S_0)-E/L.
\end{equation}
The statement permits points of $S_0$ outside $S$, overlapping arcs,
singleton arcs and full-circle arcs. For $0\le E\le L$, the loss bound is
$E/L-E^2/(2L^2)$.
\end{lemma}
\begin{outline}
Let $U$ be the union of covered elementary edges and let
$C=\sum_{a\in U}w_a$. Label each edge by the gap containing it, first
for $S_0$ and then for $S$. Two uncovered edges in the same $S_0$ gap
must share an $S$ gap. Indeed, any intervening $S$ point outside $S_0$
lies in a covering arc. That arc also contains an $S_0$ point outside
the intervening interval, so it must cross one of the two supposedly
uncovered edges. This argument uses actual arc membership and does not
assume $S_0\subseteq S$.

For either gap labeling $c$, expand the squared-gap sum as
$\sum_{a,b:\,c(a)=c(b)}w_aw_b$. The preceding implication preserves
every old pair whose two edges are uncovered. Consequently the decrease
in the squared-gap sum is at most the mass of all ordered pairs touching
$U$, namely
\[
 L^2-(L-C)^2=2LC-C^2.
\]
Divide by $2L^2$ and use $0\le C/L\le\min\{E/L,1\}$ together with
the monotonicity of $\psi$ on $[0,1]$. Finally $\psi(e)\le e\le E/L$
gives the original linear bound directly. The shared proof uses finite
partition sums and the previously verified actual circular gaps; the
real-coordinate adapter follows by telescoping edge lengths.
\end{outline}

This verifies approximation by anchored arcs while retaining the complete
source set in recovery. The quadratic saving increases the available
finite geometric credit; it does not establish economical approximations
for arbitrary tape geometries.

The anchoring premise matters: adding the isolated point $L/2$ to
$\{0\}$ lowers gap area from $1/2$ to $1/4$ despite adding zero
interval length. Small total measure without anchoring is insufficient.

For a phase $\rho$ and the first output time $s_j$ of column $j$, call
$(S_0,(I_i),E)$ an \emph{admissible cover witness} if
$\varnothing\ne S_0\subseteq\{0,\ldots,L-1\}$,
$\kappa_\rho(S_0)<\infty$, all arcs are anchored in $S_0$, their union
has length at most $E$, and
$S_\rho(s_j)\subseteq S_0\cup\bigcup_i I_i$.
Each formal witness supplies a common real subdivision containing the
arc endpoints and all raw coordinates, with the latter retaining their
values under an order-preserving embedding.
Choose any family of such witnesses and set
\begin{equation}\label{eq:cover-credit}
 H_{\rho,j}=\sup\left(\{0\}\cup
 \left\{\frac{(1-\alpha)_+^2}{2(2K+1)\kappa_\rho(S_0)}
           -\psi\!\left(\min\!\left\{\frac EL,1\right\}\right):
              \text{chosen admissible witnesses}\right\}\right).
\end{equation}
An empty witness family gives zero. The preceding lemmas show that
$0\le H_{\rho,j}\le\gaparea(S_\rho(s_j))\le1/2$, whether or not the
supremum is attained. Bounds below that subtract $E/L$ remain valid
relaxations of this stronger credit. In particular one may use the exact witness
$S_0=S_\rho(s_j)$, $E=0$. Let $\overline H$ be the arithmetic mean
of $H_{\rho,j}$ over the same legal phases and all $d$ columns.

\begin{theorem}[Theorem C: complete covers with the full paid cost]\label{thm:paid-source-cover}
\formal{SourceCovers/PaidSourceCover.lean}{SourceCovers.complete_cover_paid_deficit,SourceCovers.paid_deficit_of_gap_credit,SourceCovers.phaseCoverCredit_bounds,SourceCovers.meanCoverCredit}
Suppose every actual family outside the selected core and the $q$
exceptions has at most $a$ hops. For every integer $w\ge1$, define
\begin{align}\label{eq:full-paid-error}
 \mathsf E_q(w;J,M,R,a,\tau)
 ={}&\frac{2C_q(T-1)}{dL}+\frac{2(w-1)}L+\frac{W-b+1}{2L}\notag\\
 &+\frac\beta L\left[J+q+\frac\tau{M-R}
                    +\frac{(1+Na)C_q(T-1)}w\right].
\end{align}
Then the correctness deficit on the fixed full geometry satisfies
\begin{equation}\label{eq:paid-source-cover}
 \frac\Delta L\ge\overline H-\mathsf E_q(w;J,M,R,a,\tau).
\end{equation}
All payments in this inequality use the complete bank, with the
protected rank phase and the residual grid shift chosen compatibly.
The same bound holds with $\tau$ replaced by the hop mass of the
unprotected selected families that are nonleaves of the original core.
The finite statement does not require $T\le cL$.
\end{theorem}
\begin{outline}
At every step the formal parent relation contains the state predecessor
and every actual long work-tape hop; length-one work hops are state edges.
The selected relation uses the original core and exceptions, before phase
deletion. Pay the distinct outer lowers of those families and all actual
lowers deleted by the phase. Across every unpaid selected edge, exact-bank
source sets are nested. Every remaining live lower crosses its boundary
through a state edge or a hop of an unselected family. Full-span coloring
and the outside-family size bound therefore give at most $1+Na$ such lowers.

Index the exact list bank by its slots. This preserves both its raw value
sets and its summed variation, including coincident slots. For each legal
phase, \cref{lem:residual-source-recovery} supplies one shift whose residual
payment is at most $(1+Na)V_\rho/w$, where $V_\rho$ is that bank's actual
raw variation. The whole-record tuple at all distinct paid times is one
common transcript for every output. Local determinism and actual output
correctness imply functional recovery from this tuple and the complete raw
blocks. Thus \cref{thm:cyclic-freezing} applies to the same bank and tuple;
no global recovery or final payment estimate is assumed.

For each cover witness, telescoping in its supplied real subdivision
identifies weighted gap area with the original raw point-gap area. Apply
\cref{lem:fractional-gap,lem:anchored-erosion} to bound the witness score by
the complete frozen set's gap area. The scores, together with zero, are
bounded above by $1/2$, so their supremum satisfies the same bound without
an attainment premise. Neither the witness nor its credit depends on the
residual grid shift.

Replace the frozen areas by these credits and use
$V_\rho\le C_q(T-1)$. Average each phase's full signed inequality, keeping
credit and payment together. The mean initial payment is at most the
number of distinct selected outer lowers plus the unprotected nonleaf hop
mass divided by the actual number of legal phases. The outer count is at
most $J+q$, the hop mass at most $\tau$, and the phase count at least $M-R$.
Divide by $L$ to obtain
\cref{eq:full-paid-error,eq:paid-source-cover}. The formal proof retains the
actual variation and payment quantities before these relaxations.
\end{outline}

This theorem separates the two obligations of an impossibility proof:
establish a large geometric credit and keep the entire revelation bill
below it. Optimizing a witness is an analytic operation on the fixed
geometry, not a claim that the machine can compute it cheaply.

\begin{corollary}[Counting terminal core orbits]\label{cor:orbit-count-deficit}
\formal{SourceCovers/OrbitCountDeficit.lean}{SourceCovers.phaseCoreCover,SourceCovers.meanReciprocalCover,SourceCovers.meanReciprocalCover_bounds,SourceCovers.reciprocal_cover_paid_deficit,SourceCovers.orbit_count_paid_deficit}
Under the same hypotheses, put
\[
 \kappa_{\rho,j}=\kappa_\rho(S_\rho(s_j)),\qquad
 \mathcal R=\frac1{|\mathcal A|d}
       \sum_{\rho\in\mathcal A}\sum_{j=0}^{d-1}\frac1{\kappa_{\rho,j}},
\]
where $\mathcal A$ is the set of legal phases and $s_j$ is the first
output time of column $j$. Then
\begin{gather}
 \frac1{(2M+4)^{|\mathcal E|}}\le\mathcal R\le1,\notag\\
 \frac\Delta L\ge
 \frac{(1-\alpha)_+^2}{2(2K+1)}\mathcal R
 -\mathsf E_q(w;J,M,R,a,\tau).\label{eq:reciprocal-cover-deficit}
\end{gather}
In particular, since $|\mathcal E|\le q$,
\begin{equation}\label{eq:orbit-count-deficit}
 \frac\Delta L\ge
 \frac{(1-\alpha)_+^2}{2(2K+1)(2M+4)^q}
 -\mathsf E_q(w;J,M,R,a,\tau).
\end{equation}
\end{corollary}
\begin{outline}
The exact bank at time $s_j$ is a union of at most
$(\ell+1)^{|\mathcal E|}$ complete core orbits. Each covering root
belongs to its own orbit and hence to the bank. Regression puts it
at a time at most $s_j<T$, so it is a permitted root in the fractional
cover definition. Unit weights give
$1\le\kappa_{\rho,j}\le(2M+4)^{|\mathcal E|}<\infty$;
the lower bound uses the nonempty complete bank. Taking reciprocals
and averaging proves the bounds on $\mathcal R$.

Apply \cref{lem:fractional-gap} to each complete raw bank. Its credit
is $(1-\alpha)_+^2/[2(2K+1)\kappa_{\rho,j}]$, the score of the exact
zero-arc witness. The formal proof reuses the paid averaging argument
of \cref{thm:paid-source-cover}, with these gap lower bounds as its
credits. Factoring out the common numerator gives
\cref{eq:reciprocal-cover-deficit}; the lower bound on $\mathcal R$
then gives \cref{eq:orbit-count-deficit}.
\end{outline}

This completes the finite input to the necessary hop-tail theorem below.
Retaining the mean reciprocal can improve the finite deficit when actual
covers are cheaper than the terminal-orbit count. It preserves variation
across columns and phases without assuming economical covers for arbitrary
geometries or enlarging the uniform impossibility classes by itself.

\subsection{Two-exception cover suppliers}

Let $g$ and $e$ be respectively the newer and older exceptional maps,
ordered by decreasing owner; tied owners may be ordered either way.
If there is just one exception, take it as $e$ and let $g$ be the
identity; with no exceptions both maps are the identity.
Write
\begin{equation}\label{eq:older-closure}
 R_{e,\rho}(t)=A_\rho(t)\cup
                 \bigcup_{u\in A_\rho(t)}A_\rho(e(u)).
\end{equation}
This is the complete older-exception reachability relation. It is
reflexive, transitive and causal. With
$Z_\rho(t)=\{g(u):u\in A_\rho(t),\ g(u)\ne u\}$,
\begin{equation}\label{eq:two-exception-decomposition}
 O_\rho(t)=R_{e,\rho}(t)\cup
                       \bigcup_{z\in Z_\rho(t)}R_{e,\rho}(z).
\end{equation}
Indeed successful owners decrease strictly, so a path can use $g$
before $e$ but never conversely. Let $F_\rho(t)$ be the maximal elements
of $\{t\}\cup Z_\rho(t)$ under the order
$u\succeq_e v$ if $v\in R_{e,\rho}(u)$.
This is a partial order: a nonidentity step strictly decreases time.
The point $t$ is maximal and belongs to $F_\rho(t)$.

\begin{proposition}[Complete continuation suppliers]\label{prop:continuation-covers}
\notformalized
For two exceptions the following bounds hold for every phase and root.
\begin{enumerate}[label=(\alph*)]
\item $\kappa_\rho(S_\rho(t))\le(\ell+1)|F_\rho(t)|$.
\item Let $C_g(t)=A_\rho(t)\cup\bigcup_{u\in A_\rho(t)}A_\rho(g(u))$.
Remove from $\{e(v):v\in C_g(t),\ e(v)\ne v\}$ the points already in
$C_g(t)$, and let $s(t)$ be the number of maximal remaining points under
core reachability. Then $\kappa_\rho(S_\rho(t))\le\ell+1+s(t)$.
\item Partition $F_\rho(t)\setminus\{t\}$ into $r$ nonempty clusters
with representatives $a_1,\ldots,a_r$ and closed time hulls $I_1,\ldots,I_r$.
Fix a finite bank of time maps $(b_i)_i$ satisfying
$\{b_i(u):i\}=R_{e,\rho}(u)$ at every time $u$, as supplied by
\cref{thm:exact-source-bank}. These geometry-dependent maps are fixed
before any residual width or shift is chosen. Put
$E(t)=\sum_{v=1}^r\sum_i\Var(x\circ b_i;I_v)$.
An empty remainder has $r=E(t)=0$. Then
\begin{equation}\label{eq:cluster-gap}
 \gaparea(S_\rho(t))\ge
 \frac{(1-\alpha)_+^2}{2(2K+1)(\ell+1)(r+1)}-\frac{E(t)}L.
\end{equation}
\end{enumerate}
\end{proposition}
\begin{informalproof}
Every element of a finite partially ordered set is dominated by a maximal
one. Transitivity in \cref{eq:two-exception-decomposition} therefore replaces
its seed set by $F_\rho(t)$ without losing any endpoint. Each older closure
uses at most $\ell+1$ complete core orbits, proving (a).
For (b), first complete the core and newer exception. Any additional
endpoint is in a core orbit of an older landing seed. Core reachability
is transitive, and $C_g(t)$ is core-closed; seeds already inside it add
nothing. Maximal remaining seeds suffice, each contributing one core
orbit. The newer-only part contributes at most $\ell+1$.

For (c) use the base raw set
$S_0=x(R_{e,\rho}(t))\cup\bigcup_v x(R_{e,\rho}(a_v))$.
It has a cover by $(r+1)(\ell+1)$ complete core orbits. For each cluster
and older-bank slot, the raw range interval over $I_v$ contains its
representative value and has length at most the corresponding variation.
These intervals cover every raw source omitted from $S_0$, by (a)'s
exact closure identity. They are anchored in $S_0$ and their summed
length is at most $E(t)$. Apply
\cref{lem:fractional-gap,lem:anchored-erosion}.
\end{informalproof}

All three statements concern complete continuations after a return.
Replacing $x\circ b_i$ by the immediate input coordinate $x$ would remove
unaccounted older sources. These suppliers offer concrete sufficient
conditions for \cref{thm:paid-source-cover}; no automatic bound on their
values for all native geometries is asserted.

\section{Uniform exclusions and necessary source complexity}\label{sec:uniform}

The finite bound has constants independent of the number of selected
core families. This permits a numerical limit argument. It does not
permit an infinite source bank or a cut parameter growing with dimension.

\subsection{Original hop and variation profiles}

Fix $c,\beta>0$, $N\ge1$ and an integer $q\ge1$. In each admissible
geometry let $\mathcal K_0$ be an original $\mathrm Q_2$ core and let
$\mathcal E$ contain at most $q$ other complete families. Let $a_0$ bound
the hop count of every family outside $\mathcal K_0\cup\mathcal E$.
Sort the original core by nonincreasing hop count, with a fixed rule
for ties, and denote it by $G_1,G_2,\ldots$. Both profiles below are
zero at indices with no family. The variation profile uses this same ordering, not an ordering
by variation.

\begin{definition}[Near-full sequence and joint profile]\label{def:joint-profile}
A \emph{near-full sequence} consists of admissible fixed geometries and
nonempty correctness sets indexed by $\nu$, with common $c,\beta,N,q$,
such that
\begin{equation}\label{eq:near-full}
 d_\nu\longrightarrow\infty,\qquad
 a_{0,\nu}/L_\nu\longrightarrow0,\qquad
 \Delta_\nu/L_\nu\longrightarrow0.
\end{equation}
Its original profiles are
\[
 h_i^{(\nu)}=|G_i|/L_\nu,\qquad
 v_i^{(\nu)}=\Var(x_\nu;[p_{G_i},U_{G_i}])/L_\nu.
\]
A \emph{joint profile} is a coordinatewise subsequential limit
$(h_i,v_i)_{i\ge1}$ of both sequences of profiles. Write
\[
 \tau_R=\sum_{i>R}h_i,\qquad V_K=\sum_{i>K}v_i.
\]
\end{definition}

Such a further subsequence always exists: the family accounting in
\cref{sec:source-engine} gives
$\sum_i h_i^{(\nu)}\le c$ and $\sum_i v_i^{(\nu)}\le2c$.
Successive bounded-coordinate extractions and a diagonal subsequence
give simultaneous convergence. Passing each finite partial sum to the
limit gives $\sum_i h_i\le c$ and $\sum_i v_i\le2c$.
In particular both tails tend to zero. This does not assert that tails
of the finite profiles converge or that family topologies stabilize.

\begin{theorem}[Theorem U: necessary hop tails]\label{thm:necessary-hop-tail}
\notformalized
For every joint profile of every near-full sequence and all integers
$K\ge0$, $R\ge K$, $M\ge\max(2,R+1)$,
\begin{equation}\label{eq:necessary-hop-tail}
 \frac{\beta\tau_R}{M-R}
 \ge \frac{(1-V_K)_+^2}{4(K+1)(2M+4)^q}.
\end{equation}
Consequently, whenever $V_K<1$ and $R\ge\max(K,1)$,
\begin{equation}\label{eq:tail-floor}
 \tau_R\ge
 \frac{(1-V_K)^2}{\beta4^{q+1}(K+1)(R+2)^{q-1}}.
\end{equation}
If $q\ge2$, every such profile satisfies
$\sum_{i\ge1}i^{q-1}h_i=\infty$.
\end{theorem}
\begin{informalproof}
Fix a joint-profile subsequence and $K,R,M$. For each fixed integer
$J\ge R$, select the first $J$ original core families, distinguish its
first $K$, protect its first $R$, and rebuild the forest and ranks of
that finite selection. Absent families are simply omitted. Sorted hop
counts and their total bound imply that every omitted core family has
at most $cL_\nu/(J+1)$ hops. The finite theorem therefore applies with
\[
 a_\nu=\max\{a_{0,\nu},\lceil cL_\nu/(J+1)\rceil\},\qquad
 w_\nu=\lceil\lambda L_\nu\rceil,
\]
where $0<\lambda<1$ is also fixed before $\nu$ tends to infinity.
The finite-prefix quantities satisfy
\[
 \alpha_\nu\to V_{K,J}:=\sum_{K<i\le J}v_i,\qquad
 \tau_\nu/L_\nu\to\tau_{R,J}:=\sum_{R<i\le J}h_i.
\]
The normalized residual bill is bounded in the limit by
$\beta NC_qc^2/[\lambda(J+1)]$; the normalized movement, outer-lower
and stripe errors vanish, and $2(w_\nu-1)/L_\nu\to2\lambda$.
Using \cref{cor:orbit-count-deficit} and \cref{eq:near-full}, we obtain
\begin{equation}\label{eq:finite-prefix-limit}
 \frac{(1-V_{K,J})_+^2}{4(K+1)(2M+4)^q}
 \le\frac{\beta\tau_{R,J}}{M-R}+2\lambda
       +\frac{\beta NC_qc^2}{\lambda(J+1)}.
\end{equation}
This is an inequality for every finite $J$. Now let $J$ tend to infinity
in these numerical inequalities, and then let $\lambda$ decrease to zero.
The constant $C_q$ depends on $M,q$, not $J$, so this proves
\cref{eq:necessary-hop-tail}.

Taking $M=2R+2$ gives $2M+4=4(R+2)$ and proves
\cref{eq:tail-floor}. A finite $K$ with $V_K\le1/2$ always exists.
If $\sum_i i^{q-1}h_i$ were finite, then for $R\ge1$
\[
 (R+2)^{q-1}\tau_R
 \le2^{q-1}\sum_{i>R}i^{q-1}h_i\longrightarrow0,
\]
contrary to \cref{eq:tail-floor} for that $K$.
Every step used a finite certificate. Omitted families remained in the
residual bill until after the dimension limit; no family was discarded
by an exchange of limits.
\end{informalproof}

\begin{corollary}[One arbitrary exceptional family]\label{cor:one-exception}
\notformalized
Fix $c,\beta>0$ and $N\ge1$. There exist $\eps>0$ and $d_0$ such that
every admissible geometry with a $\mathrm Q_2$ core, at most one other
complete exceptional family, $d\ge d_0$, and $a_0\le\eps L$ satisfies
\begin{equation}\label{eq:uniform-deficit}
 \Delta\ge\eps L,\qquad
 |\mathcal C|\le2^{(1-\eps)L}.
\end{equation}
No envelope on the sorted core-family sizes is assumed.
\end{corollary}
\begin{informalproof}
If no such constants existed, choose a violating geometry for each
$\eps=1/n$ and $d_0=n$. These form a near-full sequence and have a
joint profile. Choose $K$ with $V_K\le1/2$. For $q=1$,
\cref{eq:tail-floor} bounds $\tau_R$ below by a fixed positive constant
for every $R\ge\max(K,1)$, contrary to $\tau_R\to0$.
\end{informalproof}

\begin{corollary}[Common size envelopes]\label{cor:common-envelope}
\notformalized
Fix $c,\beta,N,q$ as above and a nonnegative summable sequence $(g_i)$
such that
\begin{equation}\label{eq:envelope-test}
 \liminf_{R\to\infty}(R+2)^{q-1}\sum_{i>R}g_i=0.
\end{equation}
There exist $\eps>0,d_0$ for which \cref{eq:uniform-deficit} holds
whenever $d\ge d_0$, $a_0\le\eps L$ and $|G_i|\le Lg_i$ for all
original core families in their sorted order. In particular every
common power envelope $g_i=H/i^p$, with fixed $H>0$ and $p>q$, is covered.
\end{corollary}
\begin{informalproof}
Negating the uniform conclusion again gives a near-full sequence.
In a joint profile $h_i\le g_i$ for each $i$, so
\cref{eq:envelope-test} contradicts \cref{eq:tail-floor} after choosing
$K$ with $V_K<1$. For a power envelope,
$\sum_{i>R}H/i^p\le H R^{1-p}/(p-1)$; multiplication by
$(R+2)^{q-1}$ tends to zero when $p>q$.
\end{informalproof}

\begin{corollary}[A fixed finite skeleton]\label{cor:finite-skeleton}
\notformalized
For fixed $c,\beta,N$ and $q\ge1$, there are $\eps>0,d_0$ such that
\cref{eq:uniform-deficit} holds for every admissible geometry in which
all but at most $q$ complete families have at most $a_0$ hops,
$a_0\le\eps L$, and $d\ge d_0$.
\end{corollary}
\begin{informalproof}
Take the core empty and use $g_i=0$ in \cref{cor:common-envelope}.
The full geometry and its outside-family bound are unchanged.
\end{informalproof}

These are restrictions on a common full geometry. In the one-exception
case they settle the paid-deficit comparison without a coding conjecture.
For $q=2$, the necessary tail floor is of order $1/R$, compatible with
an inverse-square size profile; a further geometric argument is needed.
The proof supplies uniform constants by contradiction, not effective
numerical values of $\eps$ and $d_0$.

\subsection{Necessary fractional-cover complexity}

Along a fixed joint-profile subsequence, select the first $J$ families
and use $M=2R+2$ as above. For each phase and column put
$\kappa_{\nu,J,R,\rho,j}=\kappa_\rho(S_\rho(s_j))$.
All these quantities are finite and at least one. Write
\begin{equation}\label{eq:reciprocal-cover-limit}
 Q_R^*=\limsup_{J\to\infty}\limsup_{\nu\to\infty}
                \mathbb E_{\rho,j}\frac1{\kappa_{\nu,J,R,\rho,j}},
\end{equation}
where the expectation is the arithmetic mean over legal phases and
columns. This order of limits is part of the definition.

\begin{theorem}[Necessary reciprocal-cover bound]\label{thm:necessary-cover}
\notformalized
For every near-full sequence, every joint profile, any $K$ with $V_K<1$
and every $R\ge\max(K,1)$,
\begin{equation}\label{eq:necessary-cover}
 Q_R^*\le
 \frac{4\beta(K+1)\tau_R}{(1-V_K)^2(R+2)}.
\end{equation}
This holds for any fixed number $q$ of exceptions.
If $|G_i|\le HL/i^2$ and $V_K\le1/2$, then for every fixed $s\ge1$,
\begin{equation}\label{eq:cover-saturation}
 \limsup_J\limsup_\nu
 \Pr_{\rho,j}\{\kappa_{\nu,J,R,\rho,j}\le s\}
 \le \frac{16\beta(K+1)Hs}{R(R+2)}.
\end{equation}
\end{theorem}
\begin{informalproof}
Use the exact witnesses in \cref{thm:paid-source-cover}. For fixed
$J,R,K,\lambda$ the geometric coefficient tends to
$(1-V_{K,J})^2/[4(K+1)]>0$. The same error calculation as in
\cref{eq:finite-prefix-limit} bounds the limsup of the mean reciprocal
cover by that coefficient's reciprocal times
\[
 \frac{\beta\tau_{R,J}}{R+2}+2\lambda
       +\frac{\beta NC_qc^2}{\lambda(J+1)}.
\]
Take the limsup in $J$, then let $\lambda$ decrease to zero. This proves
\cref{eq:necessary-cover}. The elementary pointwise inequality
$\mathbf1_{\{\kappa\le s\}}\le s/\kappa$ and
$\tau_R\le H/R$ give \cref{eq:cover-saturation}.
\end{informalproof}

\begin{corollary}[Uniform cover criteria]\label{cor:uniform-cover}
\notformalized
Fix $c,\beta,N,q$. Each of the following defines a class with constants
$\eps>0,d_0$ for which \cref{eq:uniform-deficit} holds when
$d\ge d_0$ and $a_0\le\eps L$.
\begin{enumerate}[label=(\alph*)]
\item For some common constant $0<B<\infty$, selecting the first $J$
original core families and protecting the first $R$, for every
$J\ge R\ge1$ and $M=2R+2$, gives
$\mathbb E_{\rho,j}\kappa_{J,R,\rho,j}\le B(R+2)$.
No size envelope is required.
\item For fixed $H>0,p>1$ the original sizes satisfy $|G_i|\le HL/i^p$,
and for a common positive function $f(R)=o(R^p)$ every such finite
prefix selection satisfies $\mathbb E_{\rho,j}\kappa_{J,R,\rho,j}\le f(R)$.
\end{enumerate}
Here $\kappa_{J,R,\rho,j}$ is the fractional cover of the selected
geometry's complete source set at the first output time of column $j$.
The bounds are uniform over the class and dimension; $\eps,d_0$ may
depend on all the fixed class data. In (b), the
inverse-square case requires mean cover $o(R^2)$.
\end{corollary}
\begin{informalproof}
Suppose either uniform conclusion fails and extract a joint profile of
a near-full sequence. Jensen's inequality for $u\mapsto1/u$ gives,
respectively, $Q_R^*\ge1/[B(R+2)]$ or $Q_R^*\ge1/f(R)$.
Choose $K$ with $V_K<1$. In (a), \cref{eq:necessary-cover} then forces
$\tau_R$ to stay above a positive constant, a contradiction.
In (b), $\tau_R\le HR^{1-p}/(p-1)$ bounds its right side by
a fixed multiple of $R^{-p}$, contradicting $f(R)=o(R^p)$.
\end{informalproof}

The saturation bound also rules out a misleading average picture:
in a near-full inverse-square sequence, rare very complicated columns
cannot compensate for a positive fraction with cover complexity
$o(R^2)$, in the stated ordered limits. This is a necessary condition,
not a construction of an economical cover.

\subsection{The movement of complete older observables}

For $q=2$, use one cluster for all nonroot elements of $F_\rho(s_j)$
in \cref{prop:continuation-covers}(c), and use $E=0$ when that set is
empty or a singleton. Let $e_{\nu,J,R}$ be the mean of $E/L_\nu$
over phases and columns, with $M=2R+2$, and put
$e_R^*=\liminf_J\liminf_\nu e_{\nu,J,R}$.

\begin{corollary}[Necessary complete-observable movement]\label{cor:necessary-exposure}
\notformalized
Along every subsequence realizing a joint profile of a two-exception
near-full sequence, for every integer $K\ge0$ and $R\ge\max(K,1)$,
\begin{equation}\label{eq:necessary-exposure}
 e_R^*\ge
 \frac{(1-V_K)_+^2}{32(K+1)(R+2)}-\frac{\beta\tau_R}{R+2}.
\end{equation}
For fixed $K$ with $V_K<1$ this implies
$\liminf_{R\to\infty}(R+2)e_R^*\ge(1-V_K)^2/[32(K+1)]$.
\end{corollary}
\begin{informalproof}
The one-cluster estimate, also valid when there are no nonroot seeds,
is $(1-\alpha)_+^2/[8(K+1)(2M+4)]-E/L$.
Insert it in \cref{thm:paid-source-cover} and move the mean $E/L$ to
the other side. Take the liminf in $\nu$, then in $J$, and then let
$\lambda$ decrease to zero exactly as in
\cref{eq:finite-prefix-limit}. Substitution of $M=2R+2$ gives
\cref{eq:necessary-exposure}; $\tau_R\to0$ gives the last assertion.
\end{informalproof}

Here the moving observables include every older continuation, not just
the immediate input coordinate at an exceptional landing. This is the
quantity that a successful approximation argument would have to control.

\section{Paying shared raw support once}\label{sec:shared-support}

The preceding covers estimate source gap area without adding revelations.
Here we use a different mechanism: reveal a raw source set shared by many
outputs, pay its union once, and apply the one-exception area bound to
the sources that remain. All assertions in this section use $q=2$.

\subsection{A finite global-sharing bound}

Retain the notation $A_\rho$, $R_{e,\rho}$ and $Z_\rho(t)$ of
\cref{eq:older-closure,eq:two-exception-decomposition}. Fix
$J\ge R\ge\max(K,1)$, $M=2R+2$, and set
\begin{equation}\label{eq:sharing-constants}
 \ell=4R+7,\quad C_R=C_0(1+3C_0)^2,\quad
 B_R=\ell(\ell+1),\quad A_R=\beta NC_Rc,
 \qquad C_0=16+24M(M-1).
\end{equation}
Let $\mathcal T_{\rm out}$ be the set of all $L$ actual output-bit times,
and define
\[
 Z_\rho^{\rm all}=\bigcup_{t\in\mathcal T_{\rm out}}Z_\rho(t),\qquad
 G_\rho=\left|\max_{\succeq_e}Z_\rho^{\rm all}\right|,\qquad
 \overline G_{J,R}=\mathbb E_\rho G_\rho.
\]
An empty seed set has $G_\rho=0$. The order $\succeq_e$ is complete
older reachability, not merely temporal order. These maximal seeds are
geometric witnesses, not independent records of bounded bit capacity.
Let
\begin{equation}\label{eq:global-newer-support}
 S^{\rm new}_\rho
 =x\left(\bigcup_{z\in Z_\rho^{\rm all}}R_{e,\rho}(z)\right).
\end{equation}
Transitivity and maximality cover this entire set by the older closures
of the $G_\rho$ maximal seeds. Each has at most $B_R$ endpoints, so
$|S^{\rm new}_\rho|\le B_RG_\rho$.

\begin{proposition}[The full shared-support payment]\label{prop:finite-sharing}
\formal{SourceCovers/GlobalSharing.lean}{SourceCovers.shared_support_count_deficit,SourceCovers.shared_support_paid_deficit,SourceCovers.phase_shared_deficit,SourceCovers.sharedRawSupport,SourceCovers.sharedSeedCount,SourceCovers.exactSourceBank_shared_cover,SourceCovers.sharedRawSupport_card_le_sum,SourceCovers.sharedRawSupport_blocks_card_le,SourceCovers.exists_paidBank_shared_deficit}
With $\alpha,\tau,a$ as in \cref{sec:source-covers}, for every integer
$w\ge1$,
\begin{equation}\label{eq:finite-sharing}
 \frac\Delta L\ge
 \frac{(1-\alpha)_+^2}{2(2K+1)(\ell+1)}
 -\frac{wB_R\overline G_{J,R}}L
 -\mathsf E_2(w;J,2R+2,R,a,\tau).
\end{equation}
The additional raw revelation has at most $wB_RG_\rho$ bits for every
grid shift, so the residual-cut shift may be chosen using the original
complete two-exception bank.

More sharply, let $\overline S=\mathbb E_\rho|S^{\rm new}_\rho|$ and
$C_1=C_0(1+3C_0)$. Define $\mathsf E_{1\mid2}$ by replacing only the
first movement term $2C_R(T-1)/(dL)$ in $\mathsf E_2$ by
$2C_1(T-1)/(dL)$; its residual-cut term still uses $C_R$. Then
\[
 \frac\Delta L\ge
 \frac{(1-\alpha)_+^2}{2(2K+1)(\ell+1)}
 -\frac{w\overline S}L-\mathsf E_{1\mid2}.
\]
Neither finite bound requires a linear-runtime hypothesis. The formal
statements also allow any $M\ge2$ with $R<M$, taking $\ell=2M+3$;
the distinguished-family budget $K$ and protection budget $R$ may be
independent.
\end{proposition}
\begin{outline}
Fix a phase and select its residual shift $\theta$ for the complete
bank. In fixed raw-coordinate order reveal every bit in
\[
 U_{\rho,\theta}=
       \bigcup_{y\in S^{\rm new}_\rho}B_\theta(y).
\]
For every $\theta$, the finite-union and same-block cardinality bounds give
\[
 |U_{\rho,\theta}|\le\min\{L,w|S^{\rm new}_\rho|\}\le wB_RG_\rho.
\]
Complete older closure is unchanged when its seeds
are replaced by maximal seeds, so the raw-support count also retains
the sum of their actual closure sizes before using $B_RG_\rho$.
The source-replacement lemma applies whenever the full bank's sources
are contained in the remaining bank's sources and a common raw set.
Equal original transcripts, equal revealed raw bits, and equal remaining
source blocks then imply equal full-bank source blocks. Applying the original
recovery theorem therefore gives every output from the older bank and
one common joint transcript on the same correctness set. Its frozen
source gap has the one-exception lower bound
$(1-\alpha)_+^2/[2(2K+1)(\ell+1)]$ by
\cref{lem:fractional-gap}. Use the older bank for these sources;
its summed variation is at most $C_1(T-1)$, while the full-bank
variation in the residual-cut payment is at most $C_R(T-1)$.
\Cref{thm:cyclic-freezing} applies with total information charge
$\beta|\Gamma|+|U_{\rho,\theta}|$, without independence assumptions.
The intermediate formal bound keeps the actual union cardinality and
both actual variations. Average the complete signed inequality over
the same phases and use the existing average rank-payment bound.
This gives the sharper statement; $\overline S\le B_R\overline G_{J,R}$
and $C_1\le C_R$ give \cref{eq:finite-sharing} directly.
Revealing a seed vertex itself
would not reveal its entire source closure; the raw bits above are
what have actually been paid for.
\end{outline}

This supplies a verified finite certificate when many outputs share
complete raw support. It does not assert that such support is small
in an arbitrary tape computation; the following criteria isolate
the quantitative bounds that would suffice.

\subsection{Uniform sharing criteria}

For the uniform statements below, select the first $J$ original core
families in the hop ordering of \cref{sec:uniform}, protect the first
$R$, and rebuild the selected core's forest and ranks before choosing
a phase. All quantities indexed by $J,R$, including $\overline G_{J,R}$
and the raw covering numbers below, refer to these prefix selections.
Absent families are omitted.

Let $(\eta_J)_{J\ge1}$ be a fixed positive sequence tending to zero that
bounds the normalized size of the first omitted original core family:
$|G_{J+1}|/L\le\eta_J$, uniformly over the class. The universal choice
$\eta_J=c/(J+1)$ is available from sorted hop counts. Whenever another
bound is given, it may be replaced by its minimum with $c/(J+1)$.

\begin{theorem}[Theorem H: uniform exclusion by global sharing]\label{thm:uniform-sharing}
\formal{SourceCovers/UniformSharing.lean}{SourceCovers.sharing_uniform_deficit,SourceCovers.liminf_sharing_uniform_deficit,SourceCovers.sublinear_sharing_uniform_deficit,SourceCovers.power_sharing_uniform_deficit,SourceCovers.AdmissibleGeometry.HasSharingBounds,SourceCovers.NearFullSequence.JointProfile.prefix_sharing_limit,SourceCovers.NearFullSequence.JointProfile.sharing_tail_lower}
Fix $c,\beta,N$, a common sequence $\eta_J$ as above, and an unbounded
set $\mathcal R\subseteq\{1,2,\ldots\}$. For each $R\in\mathcal R$
fix a finite nonnegative function $g_R(J)$ on $J\ge R$ such that
\begin{equation}\label{eq:sharing-smallness}
 \liminf_{J\to\infty}\eta_Jg_R(J)=0.
\end{equation}
Consider admissible two-exception geometries for which
$\overline G_{J,R}\le g_R(J)$ for every $R\in\mathcal R$ and $J\ge R$.
These bounds and index sets are fixed independently of dimension.
There exist $\eps>0,d_0$ such that every such geometry with
$d\ge d_0$ and $a_0\le\eps L$ satisfies $\Delta\ge\eps L$, equivalently
$|\mathcal C|\le2^{(1-\eps)L}$.
The same conclusion holds for nonnegative $\eta_J$ tending to zero.

In particular it suffices, without a size envelope, that
$\liminf_J g_R(J)/(J+1)=0$ for each $R\in\mathcal R$.
Under $|G_i|\le HL/i^p$ with fixed $H\ge0$ and $p>0$, it suffices that
$\liminf_J g_R(J)/(J+1)^p=0$; this includes subquadratic sharing
for the inverse-square envelope.
\end{theorem}
\begin{outline}
Negating the conclusion gives a near-full sequence with a joint profile.
Choose a finite $K$ with $V_K<1$ and fix $R\in\mathcal R$.
For each finite $J\ge R$ use
$a=\max(a_0,\lceil\eta_J L\rceil)$ and
$w=\lceil\lambda L\rceil$, with $\lambda>0$ fixed before dimension.
Take $J\ge K$ as well. The finite sharing bound permits independent
$K,R$ and any positive integer width. Its dimension limit gives
\begin{equation}\label{eq:sharing-prefix-limit}
 \frac{(1-V_{K,J})^2}{2(2K+1)(\ell+1)}
 -\frac{\beta\tau_{R,J}}{R+2}
 \le \lambda[2+B_Rg_R(J)]+\frac{A_R\eta_J}\lambda.
\end{equation}
The finite sums $V_{K,J},\tau_{R,J}$ are as in
\cref{eq:finite-prefix-limit}, and $V_{K,J}\le V_K<1$.
Let $\mu_{J,R}$ be the left side and put $D_J=3+B_Rg_R(J)>0$.
Increase the coefficient $2$ to $3$. Whenever $\mu_{J,R}>0$, test at
\[
 \lambda=\frac{\mu_{J,R}}{2D_J}
 \quad\text{to obtain}\quad
 \mu_{J,R}^2\le4A_R\eta_JD_J.
\]
This uses no positive lower bound on $A_R$ or $\eta_J$ and no square-root
minimizer. The right side tends to zero along arbitrarily late cutoffs
provided by \cref{eq:sharing-smallness}, since $\eta_J\to0$.
The finite sums converge along those cutoffs, so their limiting
credit cannot be positive. As $\ell+1=4(R+2)$, this gives
\[
 \beta\tau_R\ge\frac{(1-V_K)^2}{8(2K+1)}.
\]
This holds for every $R\in\mathcal R$, contradicting
$\tau_R\to0$. The two special cases use respectively
$\eta_J=c/(J+1)$ and $\eta_J=H/(J+1)^p$.
Any $p>0$ makes the latter tend to zero; summability of the actual
profile was already proved and need not follow from the envelope.
The formal lower limits take values in $[0,\infty]$; zero is equivalent
to arbitrarily late values below every positive threshold.
\end{outline}

The bound $g_R$ may grow arbitrarily fast with the fixed protected
scale $R$. Its uniformity in dimension is essential: for a single
finite geometry, increasing $J$ beyond all its families proves no
uniform statement. A common unbounded set of cutoffs $J$ for each $R$
also suffices if the stated bounds and smallness hold on that set;
this is included in the formal statement. The omitted-size bound is
then needed only at those same cutoffs. For $0<p\le1$, the power-law
case is already covered by the no-envelope criterion; the broader
exponent range removes an unnecessary restriction rather than adding
a new obstruction in that range. The theorem requires these uniform
sharing bounds as hypotheses; it does not establish them for arbitrary
tape computations.

\begin{corollary}[Necessary global continuation density]\label{cor:sharing-density}
\formal{SourceCovers/SharingDensity.lean}{SourceCovers.NearFullSequence.JointProfile.seedLiminf,SourceCovers.NearFullSequence.JointProfile.necessary_sharing_density,SourceCovers.NearFullSequence.JointProfile.sharingCredit_default,SourceCovers.NearFullSequence.JointProfile.sharingCredit_eventually_pos,SourceCovers.NearFullSequence.JointProfile.sharingCost_pos,SourceCovers.NearFullSequence.JointProfile.necessary_seed_growth,SourceCovers.NearFullSequence.JointProfile.necessary_power_seed_growth,SourceCovers.NearFullSequence.JointProfile.necessary_superlinear_seed_growth,SourceCovers.NearFullSequence.JointProfile.seedLiminf_eq_top_of_zero}
Fix a near-full sequence with $c,\beta>0$ and $N\ge1$, and restrict it
to a subsequence realizing a joint profile. All dimension limits below
are along that subsequence and take values in $[0,\infty]$. For fixed $K$
with $V_K<1$ and $R\ge\max(K,1)$, define
\[
 \mu_R=\frac{(1-V_K)^2}{8(2K+1)(R+2)}-
                   \frac{\beta\tau_R}{R+2},\qquad
 g^{\rm low}_{J,R}=\liminf_{\nu\to\infty}\overline G_{J,R}^{(\nu)}.
\]
If $\mu_R>0$, then for every positive sequence $\eta_J\to0$ with
$h_{J+1}\le\eta_J$ for all sufficiently large $J$,
\begin{equation}\label{eq:sharing-density}
 \liminf_{J\to\infty}\eta_Jg^{\rm low}_{J,R}
       \ge\frac{\mu_R^2}{4A_RB_R}.
\end{equation}
The condition $\mu_R>0$ holds for all sufficiently large fixed $R$.
No bound uniform in dimension is required: it suffices to bound the
limiting hop profile. For each such $R$, without any size envelope,
the stronger conclusion is
\begin{equation}\label{eq:sharing-superlinear}
 \frac{g^{\rm low}_{J,R}}{J+1}\longrightarrow\infty.
\end{equation}
At the same fixed $R$, if $h_{J+1}\le H/(J+1)^p$ eventually,
where $H,p>0$, then
\[
 \liminf_{J\to\infty}\frac{g^{\rm low}_{J,R}}{(J+1)^p}
 \ge\frac{\mu_R^2}{4A_RB_RH}.
\]
In particular an inverse-square omitted-size bound forces positive
quadratic growth. Any common finite-dimensional size bound passes to
the limiting profile and therefore remains a sufficient hypothesis.
\end{corollary}
\begin{outline}
Fix $J$ and $\lambda>0$. Divide the finite sharing inequality by its
positive seed-count coefficient before taking the dimension lower limit.
The paid-error and width limits then give
\[
 g^{\rm low}_{J,R}\ge
 \left(\frac{\mu_{J,R}-2\lambda-A_R\eta_J/\lambda}
                  {\lambda B_R}\right)_+,
\]
where $\mu_{J,R}$ is the left side of
\cref{eq:sharing-prefix-limit}. To use only $h_{J+1}\le\eta_J$,
first replace $\eta_J$ by $\eta_J+\varepsilon$: coordinate convergence
makes this bound valid in all sufficiently large dimensions. Continuity
then removes $\varepsilon>0$. No further subsequence is needed.

For finite $g^{\rm low}_{J,R}$ and positive $\mu_{J,R}$, the shared
quadratic-tradeoff lemma gives
$\mu_{J,R}^2\le4(2+B_Rg^{\rm low}_{J,R})A_R\eta_J$.
Including nonpositive credit by its positive part yields
\[
 \eta_Jg^{\rm low}_{J,R}
 \ge\frac{(\mu_{J,R})_+^2}{4A_RB_R}-\frac{2\eta_J}{B_R}.
\]
An infinite inner lower limit satisfies this bound because $\eta_J>0$.
Now take the lower limit in $J$, using $\mu_{J,R}\to\mu_R$.
Also $\tau_R\to0$ makes $\mu_R>0$ at every sufficiently large $R$.
Extended-real scaling gives the power-growth statement.

For the stronger no-envelope conclusion, nonnegativity, monotonicity
and summability of the hop profile imply $(J+1)h_{J+1}\to0$.
The shared proof compares $(J+1)h_{J+1}$ with twice any fixed tail,
once $J$ is at least twice that tail's starting index.
Thus, for every $H>0$, the rate $H/(J+1)$ eventually bounds
$h_{J+1}$. The power-growth inequality at $p=1$ gives a lower bound
$\mu_R^2/(4A_RB_RH)$ for the normalized seed liminf. Letting $H$
decrease to zero proves \cref{eq:sharing-superlinear}.
\end{outline}

The formal density inequality also permits any $M\ge2$ with $R<M$
and independent $K,R$, and retains the positive part of the limiting
credit. At a fixed $J\ge\max(K,R)$, if $h_{J+1}=0$ and
$\mu_{J,R}>0$, its separate zero-cost case gives
$g^{\rm low}_{J,R}=\infty$.

Consequently a bound $\overline G_{J,R}\le C_R'(J+1)$, with finite
constants $C_R'$ independent of dimension, at arbitrarily large protected
scales cannot coexist with a near-full sequence. This improves the
sublinear no-envelope obstruction above.
The count concerns globally incomparable complete continuations;
it is not an information-capacity or native-gate lower bound.
No such seed-count upper bound is established for arbitrary tape computations.

\subsection{Raw metric covers instead of temporal counts}

For $\delta>0$, let $\mathcal N_{\rho,J,R}(\delta)$ be the minimum
number of closed real intervals of length at most $\delta L$ covering
$S^{\rm new}_\rho$. An empty set has covering number zero. These are
intervals in the linear raw-coordinate range, not circular arcs.

\begin{proposition}[Metric shared-support criterion]\label{prop:metric-sharing}
\formal{SourceCovers/MetricSharing.lean}{SourceCovers.metric_sharing_uniform_deficit_of_limits,SourceCovers.metric_sharing_uniform_deficit,SourceCovers.AdmissibleGeometry.prefix_metric_deficit,SourceCovers.NearFullSequence.JointProfile.prefix_metric_limit,SourceCovers.RawIntervalCover.blockUnion_card_le,SourceCovers.RawIntervalCover.two_block_factor_sharp}
In the setting of \cref{thm:uniform-sharing}, replace its bounds on
$\overline G_{J,R}$ by the following. For each $R$ in a fixed unbounded
set, there are dimension-independent numbers $\delta_{J,R}>0$
and $n_{J,R}\ge0$ with
\[
 \mathbb E_\rho\mathcal N_{\rho,J,R}(\delta_{J,R})\le n_{J,R}.
\]
Along a common unbounded set of cutoffs $J$ for that $R$, assume jointly
\begin{equation}\label{eq:metric-sharing-limits}
 \delta_{J,R}\to0,\qquad
 \delta_{J,R}n_{J,R}\to0,\qquad
 \eta_J/\delta_{J,R}\to0.
\end{equation}
Then there are $\eps>0,d_0$ such that $a_0\le\eps L$ and $d\ge d_0$
imply $\Delta\ge\eps L$ for this class.
More generally, it suffices that for each fixed $R$, every $\varepsilon>0$
and every $J_0$ admit a common allowed cutoff $J\ge J_0$ at which
$\delta_{J,R}$, $\delta_{J,R}n_{J,R}$ and $\eta_J/\delta_{J,R}$ are
all smaller than $\varepsilon$. The geometric bounds need only hold at
the allowed cutoffs; these sets remain independent of dimension.
\end{proposition}
\begin{outline}
Singleton intervals cover any finite raw set, so the set of achievable
cover cardinalities has a least element, attained by a cover. This uses
natural-number minimization; real endpoints and overlapping intervals
are permitted.

Put $w=\lceil\delta_{J,R}L\rceil>0$. A set contained in an interval of
length at most $\delta_{J,R}L\le w$ has grid labels differing by at most
one. It therefore meets at most two blocks, each containing at most $w$
raw coordinates. The union of all covered blocks has cardinality at most
\[
 \min\{L,\,2w\mathcal N_{\rho,J,R}(\delta_{J,R})\}
\]
at every shift, without a disjointness assumption on the cover. The factor
two is sharp for this finite estimate: width-one blocks and the interval
$[0,1]$ covering the two raw coordinates $\{0,1\}$ attain it.

The generalized raw-payment theorem accepts any such shift-uniform bound.
It retains the full two-exception residual bill and the older bank's
movement bound; the original per-coordinate payment is a direct
specialization. Average the new payment over the same phases and use the
existing rounded-width limit, with all cutoffs and $\delta_{J,R}$ fixed
before dimension. The right side of \cref{eq:sharing-prefix-limit} becomes
\[
 2\delta_{J,R}+2\delta_{J,R}n_{J,R}
              +A_R\eta_J/\delta_{J,R}.
\]
Simultaneous smallness at arbitrarily late cutoffs makes the limiting
credit at most zero after subtracting $\beta\tau_R/(R+2)$. The same
vanishing-tail contradiction used in \cref{thm:uniform-sharing} completes
the proof. The three joint limits in \cref{eq:metric-sharing-limits}
imply this smallness condition on their common unbounded set.
\end{outline}

The improvement from three blocks to two sharpens the finite payment
and its limiting coefficient; it does not change the vanishing conditions
or establish metric covers for an unrestricted native class.

\begin{corollary}[Sublinear raw covering exponents]\label{cor:covering-exponents}
\notformalized
The metric criterion holds in either of the following classes, with
constants uniform in dimension and cutoff and permitted to depend on
each fixed $R$ in an unbounded set.
\begin{enumerate}[label=(\alph*)]
\item Without a size envelope,
$\mathbb E_\rho\mathcal N_{\rho,J,R}(\delta)
 \le H_R\delta^{-s_R}$ for all $J\ge R$ and
$0<\delta<\delta_{0,R}$, where $H_R,\delta_{0,R}>0$ and $0\le s_R<1$.
\item Under $|G_i|\le HL/i^p$ with fixed $p>1$,
$\mathbb E_\rho\mathcal N_{\rho,J,R}(\delta)
 \le H_RJ^{u_R}\delta^{-s_R}$ in the same range, where
$0\le s_R<1$ and $0\le u_R<p(1-s_R)$.
\end{enumerate}
\end{corollary}
\begin{informalproof}
In (a) take $\eta_J=c/(J+1)$ and
$\delta_{J,R}=(J+1)^{-1/2}$. In (b) choose a fixed exponent
$u_R/(1-s_R)<v_R<p$ and set
$\delta_{J,R}=(J+1)^{-v_R}$, with
$\eta_J=\min\{c/(J+1),H/(J+1)^p\}$.
The products $\delta n$ are respectively bounded by constants times
$(J+1)^{-(1-s_R)/2}$ and $J^{u_R-v_R(1-s_R)}$; the residual ratios
$\eta_J/\delta$ also tend to zero. Eventually $\delta<\delta_{0,R}$.
Apply \cref{prop:metric-sharing}.
\end{informalproof}

These conditions concern raw-coordinate support shared by all outputs.
They can be more economical than counting separate temporal occurrences.
Supplying them for a still-open native class remains a geometric problem;
the hypotheses are not consequences of a low entropy estimate alone.

\section{Why complete continuation matters}
\label{sec:limitations}

The source-cover criteria require bounds on complete raw recovery.
Neither disjoint exceptional spans nor small variation of raw reads
transported through the exceptional maps supplies such a bound by itself.
We give one native geometry that makes this distinction explicit.  It
has a quadratic fractional cover at one root and phase, even though all
exceptional-only transported profiles vanish on the nondistinguished
core families.  The example does not satisfy a fixed inverse-square
family-size envelope and is not a counterexample to an averaged cover
criterion.

We use the strictly clipped shadows of \cref{def:clipped-shadow}, with
deleted families acting as the identity. Their ordinary orbit includes
the seed and is not truncated at other paid vertices. It is therefore
distinct from the actual dependency cone after paid lower vertices have
been removed.

For a retained core $\mathcal K$, let $B_{\mathcal K}(u)$ be the raw
labels of its ordinary orbit from time $u$.  For a nonempty raw set $S$,
define its arbitrary-seed fractional core-cover number by
\begin{equation}
 \kappa_{\mathcal K}^{\rm all}(S)
 =\inf\left\{\sum_{u=0}^{T-1}w_u:
 w_u\ge0,\quad
 \sum_{u:\,a\in B_{\mathcal K}(u)}w_u\ge1
       \text{ for every }a\in S\right\}.
 \label{eq:arbitrary-seed-cover}
\end{equation}
Allowing every computation time as a seed can only make covering easier
than requiring the seeds to lie in the original joint orbit.

\begin{proposition}[Serial exceptional spans can require a quadratic cover]
\label{prop:serial-quadratic-cover}
\notformalized
For every integer $M\ge4$ and $b\ge1$, there is a prescribed native
geometry with
\[
 d=16M,\qquad L=d^2b,
\]
eight work heads, unit moves or pauses, at most three entries per work
cell, fewer than $13L$ time vertices, and $L$ distinct increasing output
times, with the following properties.  In particular the construction
allows the logarithmic entry width $b=b(d)$.

The complete actual return families consist of a core $\mathcal K_0$
whose full spans partition into two disjoint-span classes, and two
exceptional families $F,G$ with disjoint full spans.  There are no
residual families.  A period-$M$ deletion of complete core families,
together with payment of every selected outer lower endpoint, produces
a retained core $\mathcal K$, a paid set $\Gamma$, and an output root
$t$ such that, writing $S(t)$ for its complete ordinary raw orbit and
$Q=(M-1)^2$,
\begin{equation}
 \kappa_{\mathcal K}^{\rm all}(S(t))=Q,
 \qquad
 \gaparea(S(t))\le\frac1{2Q}+\frac1{2L}.
 \label{eq:serial-cover-conclusions}
\end{equation}
The same cover equality holds if covering seeds must lie in the joint
orbit of $t$.  The paid physical dependency cone at $t$ contains $Q$
raw coordinates with circular gaps at most $\lceil L/Q\rceil$.

There is one core family $H$ such that, on the entire full span of every
other original core family, $x\circ w$ is identically zero for every
finite word $w$ in the two exceptional shadows, including the empty
word.  Here $x$ is the raw input-head position.  Thus all these
exceptional-only variation tails vanish, despite the quadratic cover.
\end{proposition}

\begin{informalproof}
Put
\[
 m=M,\quad n=8m,\quad h=4nb,
 \quad L=nh=4n^2b,
 \quad A=4L,\quad E=11L,
 \quad T=12L+h+1.
\]
All heads begin at zero.  Four heads form two copies of the following
bank construction, beginning at times $A$ and $E$.  At time $D+2ih$,
$0\le i<n/2$, head $i\bmod2$ moves up $h$, down $h$, and up $h$ from
its current base, each leg at unit speed; it waits until its next bank.
Before and after these movements it remains parked.  The base increases
by $h$ after each bank on that head.  Direct inspection of the successive
visits gives exactly the families
\begin{equation}
 \begin{split}
 C_j^D:\quad &\text{span }(D+jh,D+(j+2)h],
       \quad\text{owner }D+(j+1)h,\\
 &\text{hops }(D+(j+1)h-u,D+(j+1)h+u),\quad1\le u\le h,
 \end{split}
 \label{eq:serial-bank-families}
\end{equation}
for $D\in\{A,E\}$ and $0\le j<n$.  The two families in each bank
come from its upper and lower turn respectively.  Its interior cells
are entered on its three legs; a later bank uses a fresh spatial
interval.  The common boundary cell also has at most three entries.
Parking adds only adjacent previous visits.  Thus
\cref{eq:serial-bank-families} is the complete nonadjacent family list
on those four heads.  Even and odd $j$ give two disjoint-span classes.

A fifth head supplies $F$: it waits through $7L$, rises to $L$ by
$8L$, holds through $11L$, and descends to zero by $12L$.  Its one
complete family has
\[
 \text{span }(7L,12L],\quad\text{owner }8L,
 \quad\text{hops }(7L+a,12L-a),\quad0\le a<L.
\]
A sixth head supplies the core family $H$: it rises to $L$ by time
$L$, holds through $2L$, and descends to zero by $3L$.  Its family is
\[
 \text{span }(0,3L],\quad\text{owner }L,
 \quad\text{hops }(a,3L-a),\quad0\le a<L.
\]
There is no later ascent on either head, hence no additional lower-turn
family.  Their plateaux contribute adjacent visits.

The eighth head will supply a core bridge $B$.  Set $i=4m-1$ and
\[
 r_k=8L-(i-k)h-h/2,\qquad0\le k<m.
\]
This head waits through $5L$, rises to $L$ by $6L$, descends at unit
speed to height $m$ by $7L-m$, and then waits.  It descends by one at
each time $r_0,\ldots,r_{m-1}$ and parks at zero.  These times increase,
and $r_0>7L-m$, so the schedule is legal.  Its complete family has
owner $6L$, span $(5L,r_{m-1}]$, and hops
\[
 (5L+a,7L-a)\quad(m\le a<L),\qquad
 (5L+m-1-k,r_k)\quad(0\le k<m).
\]
In particular $B$ has $L$ actual hops.  Among bank spans it overlaps
only $C_{n-1}^{A}$: the preceding span ends exactly at $5L$.  Assign
$B$ to the even class and $H$ to that class as well.  The entire core
still has two disjoint-span classes; the early overlap graph is the
path $C_0^A,\ldots,C_{n-1}^A,B$.

The input head follows
\begin{equation}
 x(t)=
 \begin{cases}
 t,&0\le t\le L-1,\\
 2L-2-t,&L-1\le t\le2L-2,\\
 0,&t\ge2L-2.
 \end{cases}
 \label{eq:serial-raw-walk}
\end{equation}
All raw coordinates have been visited before the prescribed output
times $E,E+1,\ldots,E+L-1$.

We next specify the seventh head and the paid phase.  On the early
overlap path give $C_j^A$ rank $\min\{\max\{j,1\},n-1\}$ and $B$
rank $n-1$.  On the late path give $C_j^E$ rank
$\min\{\max\{j,1\},n-2\}$; give $H$ rank one.  Delete rank residue
zero modulo $m$.  This deletes precisely
$C_m^D,C_{2m}^D,\ldots,C_{7m}^D$ in each bank component and retains
$H,B$.  These are also the path-spine ranks with the endpoint leaves
assigned their adjacent spine rank.

Choose
\[
 t=E+(4m-1)h+h/2.
\]
Its retained late-core orbit consists of
$v_k=E+(4m-1-k)h+h/2$, $0\le k\le m-2$.
Indeed each reflection reduces the window index by one, until the
next required family is $C_{3m}^E$, which was deleted.  The actual
$F$ hop from $v_k$ lands at $r_k$.  The actual $B$ hop then lands at
$5L+a_k$, where $a_k=m-1-k$.  The last early bank family sends this
point below the bridge's outer lower endpoint:
\[
 5L+a_k\longmapsto5L-a_k<5L.
\]
Keeping the complete older continuation gives the $Q=(m-1)^2$ times
\begin{equation}
 z_{k,\ell}=A+(n-\ell)h+
 \begin{cases}
 a_k,&\ell\text{ even},\\
 h-a_k,&\ell\text{ odd},
 \end{cases}
 \quad 0\le k\le m-2,\quad1\le\ell\le m-1.
 \label{eq:serial-older-grid}
\end{equation}
They are distinct: different $\ell$ give different windows, and within
a window the offsets $a_k$ are distinct.  They are all strictly
interior to their windows, their maximum is $5L-1$, and their lowest
window is $7m+1$.  The next older family $C_{7m}^A$ was deleted, so
these are the complete older bank continuations.

Order these times as $z_1<\cdots<z_Q$ and set
\[
 u_j=2L+\lfloor jL/Q\rfloor,\qquad0\le j<Q.
\]
This ordering defines the prescribed geometry; it is not an uncharged
sorting operation in a claimed algorithm.
The seventh head waits through $u_0=2L$.  It leaves cell $j$ at time
$u_j+1$, reaches cell $j+1$, and waits until its next specified move.
Its first peak is at $u_{Q-1}+1<3L$.  It holds at height $Q$ until
$z_1-1$, then descends by one at each $z_i$ and parks.  Since both
sequences are strictly increasing, all moves are unit moves or pauses.
Its only family $G$ has
\[
 \text{span }(2L,z_Q]=(2L,5L-1],\quad
 \text{owner }u_{Q-1}+1,\quad
 \text{hops }(u_{Q-i},z_i),\quad1\le i\le Q.
\]
The lower endpoints use the last outward visits, including all waits.
The full spans of $G$ and $F$ are disjoint and can share a third
disjoint-span class.  We have now listed all $2n+4$ complete families;
the new heads have no later ascent and no omitted opposite-side
family.  All work cells have at most three entries.

Pay every actual lower endpoint of the fourteen deleted bank families
and every outer lower endpoint in the full selected family set.  The
first set has $14h$ elements.  The second has $2n+4$ elements, of
which fourteen already belong to the first.  Therefore
\begin{equation}
 |\Gamma|=14h+2n-10.
 \label{eq:serial-paid-size}
\end{equation}
Every displayed path to a $z_i$ avoids these payments: all bank offsets
are interior, all its windows avoid the deleted lower windows, and
the $F$ and $B$ hops avoid their outer lowers.  For $1\le j<Q$ the
actual $G$ and $H$ hops then give
\[
 z_{Q-j}\longmapsto u_j
       \longmapsto L-\lfloor jL/Q\rfloor.
\]
These vertices also avoid $\Gamma$.  The last $G$ hop lands at the
paid outer lower $u_0=2L$ and is discarded in the physical cone.
Together with the root's direct read zero, the surviving paths prove
that the physical cone contains
\begin{equation}
 R=\{0\}\cup
   \{L-\lfloor jL/Q\rfloor:1\le j<Q\}.
 \label{eq:serial-physical-grid}
\end{equation}
All state predecessors, other work predecessors and direct reads are
still retained; they can only add sources to this witnessed subset.
Its $Q$ circular gaps are at most $\lceil L/Q\rceil$.  Inserting
points only decreases the sum of squared gaps, so
\[
 \gaparea(S(t))\le\gaparea(R)
       \le\frac{\lceil L/Q\rceil}{2L}
       \le\frac1{2Q}+\frac1{2L}.
\]

For the exact cover assertion, the ordinary joint time orbit consists
of six groups: the $m-1$ times $v_k$, their $m-1$ images $r_k$, the
$m-1$ images $5L+a_k$, the $Q$ times $z_i$, their $Q$ strict $G$
images, and the $Q$ subsequent $H$ images.  The calculations above
show that each is reached.  They also show that the union is closed
under all retained maps: within a bank the only change is the next
listed reflection; $B$ acts on the $r_k$; $F$ acts on the $v_k$;
$G$ acts on the $z_i$; and $H$ acts on the $G$ images.  Other maps
fix these points.  At the last $G$ return the ordinary clipping rule
replaces $2L$ by $2L+1$, and $H$ then sends it to $L-1$.  Thus
\begin{equation}
 S(t)=\{0,L-1\}\cup
      \{L-\lfloor jL/Q\rfloor:1\le j<Q\}.
 \label{eq:serial-complete-raw}
\end{equation}
The extra label $L-1$ here is an ordinary clipped source; its displayed
clipped path is not asserted to survive the physical payment.

Give each of the $Q$ nonzero labels in
\cref{eq:serial-complete-raw} weight one.  A complete retained core
orbit from any computation time contains at most one of them.  To see
this, after time $3L$ every core continuation stays after $3L$, where
$x=0$.  Before time $L$ the core orbit is a singleton.  Between $L$
and $2L$, the only possible extra core image is $H$'s held owner $L$,
whose raw value is $L-2$; this is not a target because $L/Q>256$.
Between $2L$ and $3L$, the root reads zero and the sole nonidentity
core map $H$ adds at most one raw value.  The endpoints obey the same
description.  Summing the fractional covering constraints for the
$Q$ target labels therefore gives
$\sum_u w_u\ge Q$.

Conversely, use the $Q$ seeds
\[
 2L+1,\quad u_1,\ldots,u_{Q-1}.
\]
Each lies in the joint orbit and has core raw set consisting of zero
and one different nonzero target in
\cref{eq:serial-complete-raw}.  Weight one on each seed covers $S(t)$.
This proves both asserted cover equalities.

Finally, $x$ is zero throughout $[2L,T-1]$.  Both exceptional shadows
preserve that interval: their nonidentity images stay strictly above
their outer lowers, which are $2L$ and $7L$, and elsewhere they are
the identity.  Every original core span except $H$ lies in that
interval.  Hence $x\circ w$ vanishes on each such full span for every
word $w$ of exceptional shadows.  This proves the profile assertion
and completes the construction.
\end{informalproof}

The obstruction in this example is the continuation after the bridge
exit.  Many distinct intermediate times remain distinct through the
older bank, and the final core return $H$ converts them into separated
raw coordinates.  Looking only at raw reads immediately after an
exceptional hop misses this last recovery step.  In particular,
\cref{prop:serial-quadratic-cover} rules out a universal pointwise
$O(M)$ fractional core-cover bound on the stated serial and zero-tail
premises.  It also rules out a uniform pointwise lower bound
$\gaparea(S(t))\ge c/(M+1)$ with $c>0$: along this family,
$(M+1)\gaparea(S(t))\to0$.

There are two essential restrictions on that conclusion.  First, the
example ties $d=16M$ to the growing phase parameter and supplies one
root and one phase.  It does not disprove a phase or column average
whose parameters are fixed in the order required by the uniform
source-cover theorem.  Second, selecting the four complete families
$\{H,B,F,G\}$ leaves maximum family size $h=L/n=o(L)$.
Thus \cref{cor:finite-skeleton} already excludes near-full correctness
on this sequence of geometries; it does not provide a new hard native
class.  Its $2n$ equal bank families also violate every common fixed
inverse-square envelope: the $n$th largest core family has at least
$h=L/n$ hops, so a bound $h_j\le CL/j^2$ would force $C\ge n$.
These are prescribed legal geometries and source sets, not a
transposition algorithm or a proof that every reachable source carries
necessary information.  An affordable cut, an averaged cover bound,
or a different complete recovery certificate may still handle them.

\section{Movement required by a specified recovery interface}
\label{sec:movement}

A source cover must satisfy two different requirements: it must contain
enough information to recover the output, and its source positions must
have an affordable joint trajectory.  This section separates those
requirements on a familiar coded-caching example.  The code itself is
known.  The movement bounds concern its specified raw-query interface;
they do not exclude every decoder for that code.

\subsection{A finite bound independent of query order}

Write $\operatorname{dist}_L$ for distance on the circle of circumference
$L>0$.  Fix $N\ge1$ equally spaced centres
$c_i=(i+\tfrac12)L/N$, $0\le i<N$.  A query type is a $k$-element subset
of these centres, where $1\le k\le N$.

\begin{theorem}[Theorem M: movement to cover distinct query types]
\label{thm:movement}
\formal{SourceCovers/QueryMovement.lean}{SourceCovers.QueryMovement.movement_max,
SourceCovers.QueryMovement.movement, SourceCovers.QueryMovement.movement_line,
SourceCovers.QueryMovement.movement_of_charge}
Let $S_1,\ldots,S_J$ be a nonempty sequence of query types, with $F$
distinct types.  For $1\le\ell\le m$, let $y_\ell(j)$ be labelled source
positions on the circle, where $m\ge1$.  Suppose every centre in $S_j$
is within distance $\rho>0$ of some $y_\ell(j)$.  Put
\[
 V=\sum_{\ell=1}^{m}\sum_{j=1}^{J-1}
       \operatorname{dist}_L(y_\ell(j+1),y_\ell(j)),\qquad
 v=V/L,\qquad \lambda=\rho/L,
\]
\[
 s=\min\left\{N,\;m\left(\left\lfloor\frac{5\lambda N}{2}
                                  \right\rfloor+1\right)\right\},
 \qquad B=\binom{s}{k},
\]
with $\binom{s}{k}=0$ for $s<k$.  Then
\begin{equation}
 F\le (1+4v/\lambda)B.
 \label{eq:movement-finite}
\end{equation}
In particular $B>0$, and
\begin{equation}
 v\ge\frac{\lambda}{4}\left(\frac{F}{B}-1\right).
 \label{eq:movement-lower}
\end{equation}
More sharply, put
\[
 C=\sum_{j=1}^{J-1}\max_{1\le\ell\le m}
       \operatorname{dist}_L(y_\ell(j+1),y_\ell(j)).
\]
Then $C\le V$ and
\[
 F\le\left(\left\lfloor\frac{4C}{\rho}\right\rfloor+1\right)B.
\]
Both \cref{eq:movement-finite,eq:movement-lower} also hold with $v$
replaced by $C/L$.
The bounds hold for every ordering of the queries and every matching of
source positions between successive queries.  They also hold when $V$
is measured by ordinary distance on the input line.
\end{theorem}

\begin{outline}
Set $\delta=\rho/4$.  Let $a_j$ be the cumulative maximum-slot movement
before query $j$, so $a_1=0$ and $a_J=C$.  Give query $j$ the label
$\lfloor a_j/\delta\rfloor$.  There are at most
$\lfloor C/\delta\rfloor+1$ labels.  Two queries with the same label
have cumulative movement differing by less than $\delta$; the triangle
inequality therefore keeps each source within distance $\delta$ of its
position at any chosen representative query with that label.

Every centre queried with that label is consequently within $5\rho/4$
of one of the $m$ representative positions.  A circular ball of that
radius contains at most
\[
 \min\left\{N,\left\lfloor\frac{5\rho N}{2L}\right\rfloor+1\right\}
\]
of the equally spaced centres.  The formal proof lifts each covered
centre to a point of the real lattice of spacing $L/N$ inside an interval
of length $5\rho/2$.  Distinct centre indices give distinct lattice
points, whose count is at most the displayed floor bound; the count is
also at most $N$.  The union of the balls contains at most $s$ centres
and hence at most $B$ query types.  Summing over labels proves the integer
bound.  A nonempty sequence forces $B>0$.  Dropping the floor and
rearranging give the two real bounds with $C/L$; $C\le V$ gives the
original summed-movement bounds.

No step used the ordering of the types.  Any chosen matching of the $m$
positions between successive queries defines coherent labelled paths;
apply the same argument to those paths.  Finally, circle distance never
exceeds line distance, so replacing circle variation by line variation
preserves the bound.
\end{outline}

For a tape construction, this theorem charges the movement of the entire
source bank.  The sharper form requires progress even when simultaneous
slot movements are charged only by their maximum.  Relabelling the bank
slots does not remove the bound.  The shared Lean argument allows any
positive partition width $\delta$ and any nonnegative edge charge that
dominates each slot's movement; the ball radius becomes $\rho+\delta$.

\begin{corollary}[Queries spread over a column interval]
\label{cor:movement-interval}
\formal{SourceCovers/IntervalQueryMovement.lean}{SourceCovers.QueryMovement.interval_movement_max,
SourceCovers.QueryMovement.interval_movement}
Let $I_1,\ldots,I_J$ be disjoint, chronologically ordered intervals in a
finite discrete timeline.  Let $m$ labelled source trajectories have
total circle variation $V$ over that entire timeline, including the
gaps between intervals.  Assign a query type $S_j$ to each interval,
and suppose each of the $F$ distinct types occurs at least $r\ge1$
times.  For a width $w\ge0$ and every $c_i\in S_j$, suppose there is a raw point $a_{i,j}$
with
\[
 \operatorname{dist}_L(a_{i,j},c_i)\le w
\]
that is within distance $\rho>0$ of a source position at some time in
$I_j$.  The position and time may depend on $i$ and $j$.  With
$v=V/L$, $\lambda=\rho/L$, and
\[
 s'=\min\{N,\;m(\lfloor3\lambda N+2wN/L\rfloor+1)\},
 \qquad B'=\binom{s'}k,
\]
one has
\begin{align}
 F-\frac{4v}{\lambda r}&\le(1+4v/\lambda)B',
   \label{eq:movement-interval-count}\\
 v&\ge\frac{\lambda}{4}\frac{F-B'}{B'+1/r}.
   \label{eq:movement-interval-lower}
\end{align}
These bounds also hold with $v$ replaced by $C/L$, where
\[
 C=\sum_{t=0}^{T-1}\max_{1\le\ell\le m}
       \operatorname{dist}_L(y_\ell(t+1),y_\ell(t))
\]
for the full timeline $\{0,\ldots,T\}$.  More precisely,
\[
 F-\frac{4C}{\rho r}
 \le\left(\left\lfloor\frac{4C}{\rho}\right\rfloor+1\right)B'.
\]
It suffices that the intervals' internal edge sets are pairwise disjoint;
their endpoints may coincide, and no ordering is needed.  The repetition
bound $r$ may be any positive real number.
\end{corollary}

\begin{outline}
Set $\delta=\rho/4$ and charge each time step by the largest movement of
any source.  Call an interval bad if its internal charge exceeds $\delta$.
Disjoint internal edge sets give at most $C/\delta$ bad intervals.
Every type lost on removing them has all its at least $r$ occurrences
among the deleted intervals.  Counting those disjoint occurrence sets
bounds the number of lost types by $C/(\delta r)$.

Freeze the positions at the start of each remaining interval.  A
required centre is within distance $\rho+\rho/4+w$ of one of these
frozen positions: the raw point is within $w$ of its centre and $\rho$
of a source at a time whose movement from the interval start costs at
most $\delta$.  Apply the cumulative-clock count of \cref{thm:movement}
to the good interval starts, using the full timeline clock and width
$\delta$.  Each label's queried centres lie within $3\rho/2+w$ of
one representative bank, so there are at most $B'$ types per label and
at most $\lfloor C/\delta\rfloor+1$ labels.  The count applies also to
an empty set of good intervals.

Adding the lost types gives the sharper bound.  Dropping the floor and
collecting terms gives both displayed bounds with $C/L$ in place of
$v$; the denominator $B'+1/r$ is positive even when $B'=0$.
Finally $C\le V$ yields the original bounds.  If the original interval
family is empty, the result is immediate.
\end{outline}

Thus queries made at different times during an output column do not
permit us to freeze their positions for free.  Their internal movement
and the possible loss of query types both appear in the bound.
The reusable proof allows independent positive thresholds for internal
interval movement and for grouping interval starts; their sum replaces
the two occurrences of $\rho/4$ in the enlarged radius.

\subsection{A coded-caching calibration}

We use the transpose of the subset placement-delivery array in
\cite[Section IV, Theorem 2]{YanChengTangChen2017}.  Repeating its rows
and columns gives the square matrix below.  We include the elementary
decoding and entropy calculation to specify exactly which interface is
being tested.

Fix integers $N>k\ge2$ and $h\ge1$, and put
\[
 F=\binom Nk,\qquad d=NFh,\qquad L=d^2b,
 \qquad W=db,\qquad a=L/N,
\]
where $b\ge1$ is an integer.  The asymptotic application permits
$b=b(d)$.  Index the rows by $(i,u)\in\{0,\ldots,N-1\}\times
\{0,\ldots,Fh-1\}$, grouping each fixed $i$ into one contiguous bin.
Index columns by $(S,z)$, where $|S|=k$ and $0\le z<Nh$.  A column
decoder may observe all raw rows in the bins $i\in S$, together with
one common helper $G$.  The helper contains the following $b$-bit
parities, for every $(k+1)$-element set $T$ and every $u,z$:
\begin{equation}
 Z_{T,u,z}=\bigoplus_{i\in T}
              X_{(i,u),(T\setminus\{i\},z)}.
 \label{eq:subset-helper}
\end{equation}

\begin{proposition}[Exact helper and gap-area calibration]
\label{prop:subset-calibration}
\notformalized
The specified interface recovers every column on the full input cube.
For a uniform input, the helper is uniform on $g$ bits, where
\begin{equation}
 g=H(G)=L\frac{N-k}{N(k+1)}.
 \label{eq:subset-helper-size}
\end{equation}
Every helper fibre has size $2^{L-g}$.  Every entry with $i\in S$ is
absent from all helper parities and is independent of $G$ and all other
input entries.

For $P_S=\{ia:i\in S\}$ on the raw circle,
\begin{equation}
 \frac1{\binom Nk}\sum_{|S|=k}\gaparea(P_S)
   =\frac{2N-k+1}{2N(k+1)}
   =\frac gL+\frac1{2N}.
 \label{eq:subset-gap-area}
\end{equation}
Consequently, for fixed $k$ and $N\to\infty$, the queried row fraction
$k/N$ tends to zero, whereas both $g/L$ and the mean point-gap area
tend to $1/(k+1)$.
\end{proposition}

\begin{informalproof}
If $i\in S$, the desired entry is read directly.  If $i\notin S$,
take $T=S\cup\{i\}$ in \cref{eq:subset-helper}.  All the other terms
have row-bin index in $S$, so the queried rows and that parity determine
the entry.  Every nonqueried entry belongs to exactly one parity,
specified by $T=S\cup\{i\}$ and its $u,z$.  The parity classes are
disjoint, so their bits are independent and uniform.  Counting them
gives
\[
 g=b\binom N{k+1}(Fh)(Nh)
   =L\frac{N-k}{N(k+1)}.
\]
The parity map is surjective, with all fibres of size $2^{L-g}$.
Queried entries occur in none of these parity classes, proving the
independence assertion.

For the gap formula, choose $S$ uniformly and mark one of its elements
uniformly.  Starting at the marked element, the clockwise integer gaps
$D_1,\ldots,D_k$ are uniform among positive compositions of $N$.
Indeed, each composition and starting point specifies one marked set,
and $N\binom{N-1}{k-1}=k\binom Nk$.  Put $E_i=D_i-1$ and $n=N-k$.
Then $\sum_iE_i=n$.  Stars and bars gives
\[
 \mathbb E\sum_i E_i(E_i-1)
 =\frac{2k\binom{n+k-1}{k+1}}{\binom{n+k-1}{k-1}}
 =\frac{2n(n-1)}{k+1},
\]
with both sides zero when $n<2$.  The numerator follows, for example,
by taking the coefficient of $z^n$ in
$2kz^2(1-z)^{-k-2}$.  Expanding $\sum_i(E_i+1)^2$ now yields
\[
 \mathbb E\sum_iD_i^2=\frac{N(2N-k+1)}{k+1}.
\]
The geometric gaps are $aD_i$, so division by $2L^2$ proves
\cref{eq:subset-gap-area}.  The limits follow from the displayed
formulas.
\end{informalproof}

This example prevents an information-only inference from sparse queried
rows to a deficit exceeding the helper cost.  The exact recovery code
exists, and the normalized gap area approaches the same value as that
cost.  The remaining issue is its geometric realization.

\begin{corollary}[The dimension limit precedes the radius limit]
\label{cor:subset-movement}
\notformalized
Fix $k\ge2$, $m\ge1$ and $h\ge1$, and take the preceding interfaces
with $N\to\infty$.  For each $N$, allow arbitrary column order and
arbitrary coherent source matching.  Let $v_N$ be the total circle
variation divided by $L$.

If every queried bin centre is covered at the column representative
within radius $\rho=\lambda L$, then, for every fixed
$0<\lambda<2/(5m)$,
\begin{equation}
 \liminf_{N\to\infty}v_N
 \ge\frac{\lambda}{4}
      \left[\left(\frac{5m\lambda}{2}\right)^{-k}-1\right].
 \label{eq:subset-centre-limit}
\end{equation}
If instead one required raw point within $W$ of each queried bin centre
is covered within radius $\lambda L$ at some time during its column
interval, then, for every fixed $0<\lambda<1/(3m)$,
\begin{equation}
 \liminf_{N\to\infty}v_N
 \ge\frac{\lambda}{4}
      \bigl[(3m\lambda)^{-k}-1\bigr].
 \label{eq:subset-interval-limit}
\end{equation}
Both right-hand sides tend to infinity as $\lambda\downarrow0$.
Thus no fixed variation constant represents this interface with a fixed
number of source trajectories at every sufficiently small fixed radius
ratio.
\end{corollary}

\begin{informalproof}
For the first assertion, $s/N\to5m\lambda/2<1$ in
\cref{thm:movement}.  Since $k$ is fixed,
\[
 \frac{\binom{s}{k}}{\binom Nk}\longrightarrow
        (5m\lambda/2)^k.
\]
Apply \cref{eq:movement-lower}.  For the second assertion, each type
occurs $r=Nh$ times and $w=W=L/d$.  Hence
$2wN/L=2/(Fh)\to0$, $s'/N\to3m\lambda<1$, and
$\binom{s'}k/\binom Nk\to(3m\lambda)^k$.  Divide the numerator and
denominator of \cref{eq:movement-interval-lower} by $\binom Nk$ to
obtain \cref{eq:subset-interval-limit}.  Each lower bound is a positive
constant times $\lambda^{1-k}$ minus $\lambda/4$, which proves the
last assertion.
\end{informalproof}

The necessary-point premise has a precise information-theoretic use.
Choose, for each queried bin, one bit of an own-column entry in a middle
row.  Its raw coordinate is within $W$ of the bin centre.  That bit is
independent of the specified helper and every other raw bit, so a
decoder on the full cube cannot omit it from its raw dependency set.
Any geometric representation of that dependency set must cover it.
When those raw observations are represented during the column interval,
\cref{cor:subset-movement} applies.

This necessity assertion does not survive every change of helper.
Revealing all entries with $i\in S$ costs exactly $kL/N=o(L)$ extra
bits at fixed $k$ and removes these particular witnesses.  That payment
does not show that the remaining recovery is cheap; it shows why the
movement result is about the stated interface.  A native source-bank
argument may reveal additional threshold or residual information, so
its required queries must be established after those exact disclosures.
The radius ratio is also fixed before the dimension limit; the
conclusion does not interchange these limits by setting $\lambda=1/N$.

\section{Causally autonomous column groups}\label{sec:causal-groups}

The source-cover results restrict the geometry of complete dependence.
This section gives a complementary restriction on the functions computed
along that geometry. If every local update obeys a capacity-conserving
decomposition into fixed narrow groups of consecutive columns after one
helper is fixed, then a
linear-length computation must disclose almost the entire input through
that helper.  Neither the return hypotheses nor final column purity
alone supplies this causal decomposition.

\subsection{The causal hypothesis and the routing convention}

Here $b\ge1$ is any integer, $W=db$, and $L=d^2b$.  Consider a directed
acyclic network with nodes $0,\ldots,n-1$, all edges directed from smaller
to larger nodes, and all memory edges $(v,v+1)$.  Edge $e$ has a fixed
nonnegative capacity $c_e$, measured in bits.  Each input row $i$ has a
specified set $A_i$ of attachment nodes; at a node in $A_i$ the entire
raw row is visible.  Each raw bit $(i,j,t)\in[d]^2\times[b]$ has a
specified output node $o(i,j,t)$.  These nodes respect transposed output
order: the bits of column $j$ are output before those of column $j+1$,
although several columns can share a node.  These are the row-source
and output conventions of the block-network model
in~\cite[section ``Codes on time--space diagrams'']{FirstPaper2026}.

Let $X\in\{0,1\}^{L}$ have any probability law supported on correctly
transposed inputs, and let $Z$ be a deterministic function of $X$.
Partition the columns into nonempty consecutive intervals
$J_1,\ldots,J_h$, in increasing order.  Write $g_a=|J_a|$,
$g=\max_a g_a$, and $X_a=X|_{[d]\times J_a\times[b]}$.
All entropies and logarithms in this section use base~$2$.

\Needspace{17\baselineskip}
\begin{definition}[Causal group typing]\label{def:causal-group-typing}
\formal{SourceCovers/CausalGroups.lean}{SourceCovers.CausalGroups.ColumnPartition,
SourceCovers.CausalGroups.Network, SourceCovers.CausalGroups.Network.FibreTyping,
SourceCovers.CausalGroups.Network.CausalGroupTyping,
SourceCovers.CausalGroups.Network.FibreTyping.local_edge_function,
SourceCovers.CausalGroups.Network.FibreTyping.local_output_function}
The computation has a \emph{capacity-conserving causal group typing
conditional on $Z$} if the following holds for each attained $z$.
For each edge $e$ there is one tuple of fields of each type
$a\in[h]$, taking values in an alphabet of size at most
$2^{c_{e,a}(z)}$, where
\[
 c_{e,a}(z)\ge0,\qquad \sum_{a=1}^{h}c_{e,a}(z)\le c_e.
\]
Every outgoing field of type $a$ is a deterministic function only of
incoming type-$a$ fields, the raw bits of $X_a$ in rows attached at that
node, and $z$.  Each output bit in group $a$ is obtained by a function
with these same permitted arguments.  The field alphabets, allocations
and local functions are fixed on the whole fibre $Z=z$, and their
evaluation agrees there with the computation and its outputs.
\end{definition}

The Lean definition expresses local dependence by equality of the
permitted observations on the fibre. Mathlib's function-factorization
theorem supplies actual functions on all local tuples; an attained fibre
supplies default values on unattained tuples. Conversely, local functions
with the stated agreement satisfy this equality-based condition. The
formal network retains the row attachments, ordered outputs and memory
spine. A monotone surjection assigns columns to their nonempty consecutive
groups.

All input-dependent control used by a local function must consequently
be carried by its permitted fields or determined by $z$.  A shared
symbol cannot be reused as a full-capacity field for several groups
without paying the sum of those capacities.  The local functions may be
nonlinear and need only match the original computation on the fibre;
they can be extended arbitrarily to unattained local tuples.  Merely
asserting that an edge's final value depends on one group does not
establish this local causal factorization.

A \emph{partial routing} below uses the underlying undirected network.
For each raw bit it assigns mass in $[0,1]$ to walks from nodes in its
row's attachment set to its own output node.  Every edge traversal,
in either direction and with multiplicity, consumes capacity.  Bits of
unattached rows receive zero mass.  This is an analytic fractional
routing; it is neither a causal execution nor a routing at one common
rate for all bits.

For each group let $\alpha_a$ and $\omega_a$ be its first and last
output nodes, and put
\[
 I_a=\{(v,v+1):\alpha_a\le v<\omega_a\},\qquad
 W_a=db g_a.
\]
The edge intervals $I_a$ are pairwise disjoint, since the groups are
consecutive.  They may share endpoint nodes.  Define the paid volume
\begin{equation}\label{eq:group-volume}
 V=\sum_e c_e+2\sum_{a=1}^{h}W_a(\omega_a-\alpha_a)
 \le \sum_e c_e+2gdb(n-1).
\end{equation}
The factor $2$ pays both collection at $\omega_a$ and return to each
bit's original output node.

\Needspace{12\baselineskip}
\subsection{From autonomous recovery to partial routing}

\begin{theorem}[Theorem Q: paid routing from causal groups]
\label{thm:causal-group-routing}
\formal{SourceCovers/CausalRoutingBound.lean}{SourceCovers.CausalGroups.Network.causal_group_routing,
SourceCovers.CausalGroups.Network.causal_group_routing_sharp}
Assume \cref{def:causal-group-typing}.  For every attained $z$, let
\[
 \mathcal A_a(z)=
 \{x|_{[d]\times J_a\times[b]}:x\in\operatorname{supp}(X\mid Z=z)\},
 \qquad
 R_{\mathrm{sup}}=
 \sum_z\Pr[Z=z]\sum_{a=1}^{h}\log_2|\mathcal A_a(z)|.
\]
There is a partial routing of total mass $R_{\mathrm{sup}}$, with
capacity on edge $e$ at most
\begin{equation}\label{eq:group-augmented-capacity}
 c'_e=c_e+2\sum_{a:e\in I_a}W_a.
\end{equation}
For an integer $\ell\ge1$, let $S_\ell$ be the number of raw bits
whose row has an attachment and whose distance from its nearest row
attachment to its original output node is less than $\ell$.  Then
\begin{equation}\label{eq:causal-group-finite}
 \sum_{a=1}^{h}H(X_a\mid Z)
 \le R_{\mathrm{sup}}
 \le S_\ell+\frac{V}{\ell}.
\end{equation}
In particular, writing $\Delta=L-H(X)$ and
\[
 \operatorname{TC}(X_1;\ldots;X_h\mid Z)
   =\sum_a H(X_a\mid Z)-H(X\mid Z),
\]
one has
\begin{equation}\label{eq:causal-group-helper}
 \Delta+H(Z)\ge
 L+\operatorname{TC}(X_1;\ldots;X_h\mid Z)
   -S_\ell-\frac{V}{\ell}.
\end{equation}
No independence of the input rows, groups, or helper fibres is assumed.
\end{theorem}

\begin{outline}
Represent the input law by weights on its exact support $C$; all these
weights are positive. The helper is indexed by its attained values.
Fix $a,z$.  Causal induction shows that evaluating all fields of type
$a$ depends only on $X_a$ and $z$.  Indeed, this is true of the locally
visible raw data, and it is preserved at each node by the local typing
rule.  Thus a projected input $y\in\mathcal A_a(z)$ defines a code
recovering every bit of $y$, independently of which full input extending
$y$ was used to define it.  This argument evaluates one projection at a
time; it makes no product assumption about the conditional support.

Add a forward collector along each edge of $I_a$. At edge $(v,v+1)$ it
carries precisely the group bits whose original output node is at most
$v$, so it needs only their number of bits in capacity, at most $W_a$.
At node $\omega_a$ this prefix together with the outputs at that node
determines all of $y$. These collector fields still obey group autonomy.
Their exact total capacity is
\[
 \sum_{\xi\in[d]\times J_a\times[b]}(\omega_a-o(\xi))
 \le W_a(\omega_a-\alpha_a),
\]
including zero when both endpoints coincide. The equality follows by
counting, for each output bit, the collector edges that carry it.

Add a supersource $s$ and a source vertex $u_i$ for each row.  Give
$(s,u_i)$ capacity $bg_a$, and give every arc from $u_i$ to a node of
$A_i$ capacity $W_a$.  The sink is $\omega_a$.  Consider any
$s$--$\omega_a$ cut.  If it crosses one of the latter attachment arcs,
its capacity is at least $W_a\ge\log_2|\mathcal A_a(z)|$. Otherwise let
$U$ be its sink-side physical nodes and let $B$ be its sink-side row
vertices.  Every row attached in $U$ belongs to $B$.

For an input $y\in\mathcal A_a(z)$, record its $bg_a$ bits in each
row of $B$ and all type-$a$ fields and collector values entering $U$.
This tuple, with $z$ fixed, determines the computation inside $U$ in
node order.  It supplies every raw argument visible there and every
parent value entering from outside, including any edge that re-enters
$U$ after an earlier departure.  It therefore determines $y$ at the
collector sink.  The tuple map is injective, giving
\[
 \log_2|\mathcal A_a(z)|
 \le bg_a|B|+\text{capacity entering }U.
\]
The right side is the cut capacity. The directed max-flow/min-cut theorem,
whose original undirected form is due to Ford and Fulkerson
\cite{FordFulkerson1956}, supplies a flow of value
$\log_2|\mathcal A_a(z)|$, truncated to that value if necessary.
For completeness, the directed version with real capacities follows by
maximizing flow value over the compact feasible-flow polytope. A residual
walk to the sink gives a signed change on the original edges: divergences
cancel at its internal vertices, and a sufficiently small positive step
respects every capacity bound. Such a walk would increase the maximum.
Without one, the residual reachable set gives a cut of the same value.
Row $i$ supplies at most $bg_a$, and an unattached row supplies zero.

Extract weighted source-to-sink paths carrying the full flow value.
A positive remaining source value forces a path of positive-flow edges,
by conservation across the reachable cut. Subtract the largest permissible
path weight, capped by the remaining source value. This either finishes
or exhausts a positive edge, so at most one path per originally positive
edge is needed. Every traversal is charged with multiplicity.
Remove the paths' artificial source and attachment arcs.
Split each path's weight equally among the $bg_a$ original bits in its row.
Each bit receives that row's total path mass divided by $bg_a$, hence at
most one unit. This preserves every retained edge's load. All these bits
have the same permitted row attachments; no particular bits need to be
identified as entropy carriers.
For each allocated bit append the memory path backwards from
$\omega_a$ to that bit's original output node. This return crosses
$(v,v+1)$ exactly when the bit's output is at most $v<\omega_a$.
The unit bit caps therefore bound the return load by the same prefix
size used by the collector. Their combined charge is at most twice
that prefix size, hence at most $2W_a$ per edge of $I_a$; its total is
at most $2\sum_\xi(\omega_a-o(\xi))$ over the group's bits. Collapsing
collector copies onto original memory edges adds their traversal loads.

Sum over groups.  On an original edge the group flows use at most
$\sum_a c_{e,a}(z)\le c_e$.  The collector and return loads obey
\cref{eq:group-augmented-capacity}.  Average these routings over $z$.
Their common capacity constraints and per-bit mass bounds are linear,
so the average is feasible and has total mass $R_{\mathrm{sup}}$.
Also $H(X_a\mid Z=z)\le\log_2|\mathcal A_a(z)|$, proving the
first inequality in \cref{eq:causal-group-finite}.

At most $S_\ell$ routed mass belongs to short-distance bits, since each
bit has mass at most one.  Every unit of the remaining mass traverses
at least $\ell$ physical edges. Formally, ordinary undirected graph
distance is bounded by each physical walk's length, and total traversal
load equals that length. Taking the infimum over actual row attachments
gives the nearest-attachment distance, with infinity for an unattached row.
Consequently
$\ell(R_{\mathrm{sup}}-S_\ell)\le V$, proving the second inequality.
Unattached rows contribute no routed mass and are not counted as
zero-distance sources.  Finally, because $Z$ is determined by $X$,
\[
 \sum_aH(X_a\mid Z)
 =L-\Delta-H(Z)+\operatorname{TC}(X_1;\ldots;X_h\mid Z).
\]
Substitution proves \cref{eq:causal-group-helper}.
The conditional coordinate-cover inequality also verifies that the total
correlation is nonnegative. The formal theorem retains the sharper prefix
charges above; the stated interval-budget version follows directly.
Finally, disjoint collector links inject into the $n-1$ memory slots,
giving the last inequality in \cref{eq:group-volume}.
\end{outline}

The theorem supplies an explicit conversion from a restricted causal
code to partial routing.  Its extra communication is fully charged.
It does not strengthen the general routing obstruction, and it does
not convert arbitrary causal codes to routings.

\subsection{Application to linear-length tape diagrams}

For clarity we state the precise short-distance estimate needed for
partial, possibly unequal, commodity masses.  The summed-distance bound
alone would not suffice.

\Needspace{10\baselineskip}
\begin{lemma}[Attached short-distance count]\label{lem:attached-short-count}
\formal{SourceCovers/NativeShortCount.lean}{SourceCovers.attachedShortBits, SourceCovers.card_attachedShortBits_le_local, SourceCovers.card_attachedShortBits_le_profile, SourceCovers.card_attachedShortBits_le, SourceCovers.attachedShortBits_zero_radius, SourceCovers.attachedShortBits_zero_block}
Let a $k$-work-tape diagram have block length $r\ge0$, $n$ block
nodes, and raw length $L=d^2b$. Use its original row attachments and
transposed-bit output nodes. Write $R(a)$ for the rows attached at node
$a$, $B_t(a)$ for its closed undirected radius-$t$ ball, and $m_{i,s}$
for the number of bits of raw row $i$ assigned to output node $s$.
For every integer $\ell\ge1$,
\begin{equation}\label{eq:attached-short-local}
 S_\ell\le
 \sum_{a<n}\sum_{i\in R(a)}\sum_{s\in B_{\ell-1}(a)}m_{i,s}.
\end{equation}
This bound requires neither diagram validity nor an upper bound on $r$.
If the diagram is valid for output length $L$ and $r\le db$, then
\begin{align}
 S_\ell&\le b\sum_{a<n}|R(a)|\,|B_{\ell-1}(a)|,
     \label{eq:attached-short-profile}\\
 S_\ell&\le\min\{L,\,(2n-1)_+b(2k+3)^{\ell-1}\}.
     \label{eq:attached-short-count}
\end{align}
Also $S_0=0$, and $S_\ell=0$ for $r=0$ under the library's zero-block
convention. The Lean inequalities include these cases, interpreting
natural subtraction $\ell-1$ as $\max\{\ell-1,0\}$ and
$(2n-1)_+=\max\{2n-1,0\}$.
\end{lemma}

\begin{outline}
The formal definition excludes unattached rows, whose inherited default
distance is zero. If an attachment exists, the block network is nonempty,
its chosen output node lies in that network, and the memory spine
connects it to the attachment. These facts require no validity
hypothesis. The bridge lemma therefore identifies a short bit with a
row attachment $a$ whose radius-$(\ell-1)$ ball contains its output
node. Taking the union over $a$, its attached rows, and its ball gives
\cref{eq:attached-short-local} by Mathlib's finite-union bound.

For a valid diagram and $1\le r\le db$, a block's output positions
lie in an interval of length at most $db$. Two distinct bits of one raw
row with the same within-entry offset have transposed positions separated
by a nonzero multiple of $db$. They cannot both lie in this interval.
The within-entry offset thus injects the bits counted by $m_{i,s}$ into
$\{0,\ldots,b-1\}$, proving $m_{i,s}\le b$ and
\cref{eq:attached-short-profile}. This sharpens the first paper's
$2b$ estimate, which separately counts the portions in two output
columns~\cite[section ``Codes on time--space diagrams'']{FirstPaper2026}.

A block attaches at most two raw rows, and its underlying graph has
degree at most $2k+2$. Reuse the first paper's
\lean{card_rowsAt_le} and \lean{card_ball_le} from
\leanfile{Transpose/Tape/CECount.lean} to obtain
$|R(a)|\le2$ and $|B_t(a)|\le(2k+3)^t$.
The first block can touch only row zero: the input head starts at zero
and has fewer than $r\le db$ steps in that block. Consequently
$\sum_{a<n}|R(a)|\le(2n-1)_+$. Together with $S_\ell\le L$,
this gives \cref{eq:attached-short-count}. The zero-radius set is empty by its
strict distance test; for zero block length there are no nodes or
attachments.
\end{outline}

The profile bound retains actual row and ball counts for later use,
without imposing a common routing rate on the counted bits. The sharper
coefficient and first-block correction improve the finite estimate;
they do not change the asymptotic obstruction below. No optimality of the
graph-ball estimate or the union bound is claimed.

\Needspace{8\baselineskip}
\begin{corollary}[Narrow causal groups require an almost complete helper]
\label{cor:causal-group-obstruction}
\notformalized
Fix $k\ge0$, an integer $G\ge1$, and $c>0$.  Along any sequence
$d\to\infty$, let $b\ge1$, $L=d^2b$, and
$r=\lceil db/G\rceil$.  Suppose valid $k$-work-tape diagrams have
length $T\le cL$, and their code capacities satisfy the block-profile
volume bound
\[
 \sum_e c_e\le4knrG,
 \qquad n=\lceil T/r\rceil.
\]
Let $X$ be any law on correctly transposed inputs and $Z$ a
raw-input-determined helper admitting causal group typing into fixed
consecutive column groups of maximum size $g$.  Then, with constants
depending only on $k,G,c$,
\begin{equation}\label{eq:causal-group-asymptotic}
 \sum_aH(X_a\mid Z)
 =O\!\left(\frac{L}{\sqrt d}+\frac{(1+g)L}{\log d}\right).
\end{equation}
In particular, if $g=o(\log d)$, then
\[
 \Delta+H(Z)\ge L-o(L).
\]
Thus $\Delta=o(L)$ and $H(Z)=o(L)$ cannot both hold.  The conclusion
applies in particular to the series' standard width
$b=b(d)=\max\{1,\lceil\log_2d\rceil\}$.
\end{corollary}

\begin{informalproof}
Since $r\ge db/G$, one has $n\le cGd+1=O(d)$.  Also
$nr\le T+r$ and $rG\le db+G$, so
$\sum_e c_e=O(L)$.  Equation~\eqref{eq:group-volume} therefore gives
$V=O((1+g)L)$.
For sufficiently large $d$, take
$\ell=\lfloor\tfrac12\log_{2k+3}d\rfloor\ge1$.
Then $\ell$ is of order $\log d$, and
\cref{lem:attached-short-count} gives
$S_\ell=O(db\sqrt d)=O(L/\sqrt d)$.
Apply \cref{thm:causal-group-routing} to obtain
\cref{eq:causal-group-asymptotic}.  Finally,
$L-\Delta-H(Z)=H(X\mid Z)\le\sum_aH(X_a\mid Z)$.
The right side is $o(L)$ when $g=o(\log d)$.
\end{informalproof}

At the standard width, the block-profile volume hypothesis is supplied
by the Lean-verified \lean{sum_cap_le} in~\cite{FirstPaper2026}.
The finite theorem also
permits other profiles: a separate $O(L)$ total capacity allowance
leaves the asymptotic conclusion unchanged.

For the construction program, the corollary excludes a completion that
preserves one fixed narrow consecutive grouping at every local update
after a sublinear-entropy joint helper is fixed.  Arbitrarily many
updates and arbitrary nonlinear operations inside each group remain
allowed.  It does not exclude final column purity, which every correct
transposer already has, nor a succession of small groups that migrate
and cumulatively mix all columns.  For nonconsecutive groups the exact
collector cost in \cref{eq:group-volume} is still available, but its
$2gdb(n-1)$ upper bound need not hold.  Finally, none of the geometric
hypotheses elsewhere in this paper is asserted to produce this causal
typing with a cheap helper; $Z=X$ always gives a vacuous expensive one.

\section{What remains for a transposition algorithm or impossibility proof}
\label{sec:roadmap}

The main conclusion is a uniform deficit for specified full return
geometries. It is stronger than a statement that one proposed code has
failed: every deterministic local computation on any geometry in the
stated class is covered, with no linearity assumption. Its scope is also
explicit. A bounded full-span partition, a Q2 core, and a bounded number
of exceptional families are additional geometric hypotheses. They have
not been derived for an arbitrary linear-time transposer.

\subsection{The nearest unresolved geometric case}

For one exceptional family, \cref{cor:one-exception} completes the
deficit-versus-payment argument. With two exceptional families, a
limiting inverse-square hop profile is compatible with
\cref{thm:necessary-hop-tail}. The next concrete problem is to determine
whether the remaining complete continuation can satisfy all the
necessary conditions simultaneously.

There are three precise ways to exclude a common native class in this
case. First, bound the phase-and-column mean fractional cover by a
common function $f(R)=o(R^2)$, uniformly in dimension and every finite
first-$J$ selection with $J\ge R$. Second, bound the mean number of
globally maximal newer seeds, at each fixed protected scale $R$ in a
common unbounded set, by a function $g_R(J)$ with
$\liminf_J g_R(J)/(J+1)^2=0$. Third, give the raw metric covers with
the three joint limits in \cref{eq:metric-sharing-limits}.
These are respectively the criteria in
\cref{cor:uniform-cover,thm:uniform-sharing,prop:metric-sharing}.
They are alternatives; none is currently supplied for every admissible
two-exception geometry with an inverse-square envelope.

A failed attempt should be checked against the exact hypothesis it
tests. \Cref{prop:serial-quadratic-cover} defeats a universal pointwise
linear cover estimate based only on serial exceptional spans and
exceptional-only variation tails. It does not defeat a phase average,
a common inverse-square envelope, or shared raw revelation across all
outputs. It also belongs to a class already excluded by the finite
skeleton corollary. Using it as evidence against all three uniform
criteria would therefore draw a conclusion the example does not support.

Conversely, a candidate near-full sequence in the inverse-square class
must meet the necessary reciprocal-cover, complete-observable movement,
and global-continuation-density bounds of
\cref{thm:necessary-cover,cor:necessary-exposure,cor:sharing-density},
with all limits taken along the same joint profile. Producing such a
geometry would test the sharpness of these certificates. It would not
by itself supply a correct computation on the geometry.
Even without a size envelope, \cref{cor:sharing-density} forces
$g^{\rm low}_{J,R}/(J+1)\to\infty$ at every sufficiently large fixed
$R$. A dimension-uniform linear seed bound would therefore already
exclude a near-full sequence in that class.

\subsection{What a constructive advance must compute}

A recovery code becomes a transposition algorithm only after its
encoding, updates, reads and output order have been implemented within
the native time bound. The complete-source certificate charges one
such implementation; it does not construct one from an abstract code.
The coded-caching calibration makes this distinction concrete. Its
algebraic decoder exists, but its specified raw-query interface cannot
be represented by a fixed number of trajectories with uniformly bounded
normalized movement at all small fixed radius ratios. Changing that
interface or helper may evade the argument and must be analysed anew.

The causal-group result identifies another restricted shortcut that
cannot work: after a cheap helper is fixed, a locally autonomous
computation for fixed groups of $o(\log d)$ adjacent columns cannot
retain nearly all the input information. A successful algorithm can
avoid that premise by coupling groups during computation. The theorem
does not price arbitrary such coupling. The constructive target is an
explicit coupled update rule whose total native execution cost can be
bounded, including the cost of changing where its encoded information
is available. Final-output separation alone is not the forbidden
causal typing.

\subsection{What would decide (T)}

An affirmative answer requires one fixed finite-tape, finite-alphabet
machine, constants $c,d_0$, and a proof that it correctly transposes
every input of every dimension $d\ge d_0$ in at most $cd^2b(d)$ steps.
All preprocessing and recovery must be included. A family of
dimension-dependent machines, a source relation, or a storage code
without such an implementation does not meet this specification.

A negative answer must exclude every such fixed machine and constant.
One possible route is a reduction from an arbitrary candidate to
correctness sets and paid certificates with a strictly positive net
deficit. Such a reduction must account for the cost of selecting a
geometry in the adaptive case, and it must address geometries outside
the span and core hypotheses used here. Alternatively, a different
certificate may handle those geometries directly. The depth reductions
of the fourth and eighth papers and the complete-recovery restrictions
of the fifth provide complementary tools
\cite{FourthPaper2026,EighthPaper2026,DirectedRecovery2026}; this paper
does not prove that their hypotheses cover every remaining machine.

The formalization verifies the finite certificates and their stated
uniform consequences. It does not promote a geometric criterion into
a universal machine property. The following appendices record that
boundary, the dependencies and the verification commands.

\appendix
\section{Specification for subsequent formalization}
\label{sec:formalization}

\Cref{def:complete-families,lem:family-accounting,def:clipped-shadow,lem:shadow-variation,lem:protected-phase,thm:exact-source-bank,lem:interval-rounding,lem:attached-short-count}
are Lean-verified;
their declarations and axiom commands accompany the paper. The remaining results retain
complete written proofs. The index in
\cref{sec:theorem-index} gives a stable label for every numbered result
and definition; it is generated from the manuscript, rather than from
a separate list of proposed theorems. Each result's status is explicit.
The statements include the hypotheses to preserve when creating Lean
declarations. Further formalizations must identify their declarations
and disclose any change of scope before calling them authoritative.

\subsection{Dependency order}

\begin{enumerate}
\item \emph{The model and complete-source engine.}
Actual last-visit hops, full first-maximum families, the half-open span
partition and family accounting are now formalized by reusing the eighth
paper and finite incidence counting. Strict clipped shadows and their
variation bounds, including successful words and disjoint packing, are
also formalized, as are protected nonleaf phases and the exact packed bank,
including closure across selected unpaid actual edges. Formalize paid-root
semantics and correctness on one fixed geometry, then residual recovery. The any-argument-path bound now covers metric-valued paths;
it is not restricted to the original input-head timeline.

\item \emph{Finite deficit.}
Formalize circular gap area, a common cyclic ordering of all raw
coordinates, stripe discrepancy, and entropy recovery from a deterministic
helper. \Cref{thm:cyclic-freezing} permits a helper bounded in entropy;
it does not infer a support-cardinality bound from entropy. The fractional
interval-cover lemma is already formalized independently of this base.
Use it to prove the fractional gap bound, formalize anchored erosion,
then prove \cref{thm:paid-source-cover}. The complete bank and the
phase used for its payment must be the same throughout.

\item \emph{Uniform consequences.}
Extract only the numerical joint profile. Formalize the order
``finite parameters, dimension limit, cutoff limit, radius limit''
in \cref{thm:necessary-hop-tail}. Derive one-exception, common-envelope,
finite-skeleton and necessary-cover results from it and the finite
certificate. Do not replace those ordered limits by a single growing
finite certificate.

\item \emph{Two-exception refinements.}
Prove the exact older-closure decomposition and its maximal-seed
reductions before the exposure and shared-support results. The global
seed set ranges over all output-bit times. Its raw support is revealed
once, at a cost valid for every grid shift, and all density limits use
the subsequence defining the same joint profile. The metric criteria
then follow by replacing temporal counts with raw interval covers.

\item \emph{Independent construction tests.}
Formalize the serial native geometry as an exact finite construction,
including all its return families, payments and arbitrary-seed cover
lower bound. Separately prove the movement theorem and its specified
query-interface applications. Finally formalize the causal-group
support-to-flow argument, followed by its native metric application.
The attached short-distance count for that application is already
formalized, including its bridge to the nearest-attachment distance.
These three branches do not serve as premises of the main uniform
exclusion theorem.
\end{enumerate}

\subsection{Reuse and statement fidelity}

The first, fourth, fifth and eighth papers contain relevant tape,
dependency, entropy and owner-order libraries. The native short-pair
count in \cref{sec:causal-groups} already imports the first paper's
row and ball counting declarations, and strengthens its output-block count.
The imported results are read from their original sources through
the existing build adapter in the second paper; no copy is introduced.
Family accounting imports the eighth paper's canonical family library
and strengthens the bound by retaining the owner's two distances to
the span endpoints. A shared weighted incidence identity supplies both
span-length and variation accounting; the original Q2 estimates follow
by bounding the overlap by two.
The clipped-shadow module reuses those actual families and the shared
monotone-patch library. Threshold crossing counts, Mathlib's finite-path
triangle inequality and order-convex preimage lemmas give the variation
bounds without a separate induction for each shadow word or packing.
The proof retains exact crossing counts and handles metric-valued paths.
The phase lemma reuses Mathlib's finite paths, cycle rotation and cardinality
bounds. Shared interval geometry, periodic phase counting and finite-orbit
modules keep the three arguments separate. The stronger cut deletes only
original nonleaves and charges their actual lower supports. The complete
orbit equality explicitly accounts for inactive actions.
The packed-bank proof reuses the patch bound and Mathlib's uniqueness of
paths in a forest. A common set of rank residues labels every spine
component injectively. Unordered endpoint pairs index the packed maps;
list composition handles the leaf sandwich and exceptional families.
The exact orbit proof permits the representing word of a packed slot to
vary with the argument. All originally selected unpaid actual edges are
covered, since the phase pays every lower of each deleted family.
The interval-cover lemma imports the shared library's greedy cover and
packing certificate. It permits an empty target and arbitrary order-convex
sets in a linear order. Its fractional bound uses Mathlib's finite-sum
and order tools; no linear-programming development is needed.
For the remaining results, inspect exact types and dependency pins
before choosing imports. Use sibling packages
through relative paths and shared lemmas through the repository's
common library; do not duplicate the earlier developments or change
their frozen theorem statements.

The following distinctions are part of the specification: a physical
head is not a return family; an endpoint witness is not a record with
free capacity; a source overcover is not a set of independent bits;
and a numerical subsequential profile is not a limiting geometry.
The movement theorem quantifies over order and matching but retains
its specified query interface. The causal theorem assumes local
functional separation on every helper fibre, not merely separated
final outputs. The serial example has a prescribed geometry and one
root and phase, not an implemented transposer or an averaged
counterexample.

After each formalization stage, compare the entire English statement
with the Lean type, including inherited model assumptions and the
order of quantifiers. The axiom commands and declaration index are
generated from the actual formal citations and must be regenerated
when those citations change. Neither the statement index nor a
successful PDF build constitutes formal verification.

% Generated by tools/statement_index.py; do not edit by hand.
\section{Statement index}\label{sec:theorem-index}
The manuscript contains 32 numbered results and 6 numbered
definitions. Lean-verified results: 6; results with complete written
proofs only: 26. A star marks a Lean-verified statement.
Stable labels below identify the statements; they are not Lean declaration names.
Source files and statement fingerprints are listed in \path{FORMALIZATION-INDEX.md}.
\begingroup\small
\begin{longtable}{@{}p{.14\textwidth}p{.39\textwidth}p{.41\textwidth}@{}}
\toprule Reference & Statement & Stable label \\
\midrule\endfirsthead
\toprule Reference & Statement & Stable label \\
\midrule\endhead
\bottomrule\endfoot
\cref{def:complete-families}${}^{*}$ & Complete families and the two-class condition & \leanname{def:complete-families} \\[2pt]
\cref{def:admissible-geometry} & Admissible return geometry & \leanname{def:admissible-geometry} \\[2pt]
\cref{lem:family-accounting}${}^{*}$ & Family accounting & \leanname{lem:family-accounting} \\[2pt]
\cref{def:clipped-shadow}${}^{*}$ & Paid clipping and selected orbits & \leanname{def:clipped-shadow} \\[2pt]
\cref{lem:shadow-variation}${}^{*}$ & Variation of a shadow or successful path & \leanname{lem:shadow-variation} \\[2pt]
\cref{lem:protected-phase}${}^{*}$ & The Q2 overlap forest and protected phases & \leanname{lem:protected-phase} \\[2pt]
\cref{thm:exact-source-bank}${}^{*}$ & Theorem S: an exact source bank & \leanname{thm:exact-source-bank} \\[2pt]
\cref{lem:residual-source-recovery} & Paid residual-source recovery & \leanname{lem:residual-source-recovery} \\[2pt]
\cref{thm:cyclic-freezing} & Theorem F: cyclic freezing & \leanname{thm:cyclic-freezing} \\[2pt]
\cref{def:fractional-core-cover} & Fractional complete-core cover & \leanname{def:fractional-core-cover} \\[2pt]
\cref{lem:interval-rounding}${}^{*}$ & Fractional interval covers & \leanname{lem:interval-rounding} \\[2pt]
\cref{lem:fractional-gap} & A fractional raw-source gap bound & \leanname{lem:fractional-gap} \\[2pt]
\cref{lem:anchored-erosion} & Anchored erosion of circular gaps & \leanname{lem:anchored-erosion} \\[2pt]
\cref{thm:paid-source-cover} & Theorem C: complete covers with the full paid cost & \leanname{thm:paid-source-cover} \\[2pt]
\cref{cor:orbit-count-deficit} & Counting terminal core orbits & \leanname{cor:orbit-count-deficit} \\[2pt]
\cref{prop:continuation-covers} & Complete continuation suppliers & \leanname{prop:continuation-covers} \\[2pt]
\cref{def:joint-profile} & Near-full sequence and joint profile & \leanname{def:joint-profile} \\[2pt]
\cref{thm:necessary-hop-tail} & Theorem U: necessary hop tails & \leanname{thm:necessary-hop-tail} \\[2pt]
\cref{cor:one-exception} & One arbitrary exceptional family & \leanname{cor:one-exception} \\[2pt]
\cref{cor:common-envelope} & Common size envelopes & \leanname{cor:common-envelope} \\[2pt]
\cref{cor:finite-skeleton} & A fixed finite skeleton & \leanname{cor:finite-skeleton} \\[2pt]
\cref{thm:necessary-cover} & Necessary reciprocal-cover bound & \leanname{thm:necessary-cover} \\[2pt]
\cref{cor:uniform-cover} & Uniform cover criteria & \leanname{cor:uniform-cover} \\[2pt]
\cref{cor:necessary-exposure} & Necessary complete-observable movement & \leanname{cor:necessary-exposure} \\[2pt]
\cref{prop:finite-sharing} & The full shared-support payment & \leanname{prop:finite-sharing} \\[2pt]
\cref{thm:uniform-sharing} & Theorem H: uniform exclusion by global sharing & \leanname{thm:uniform-sharing} \\[2pt]
\cref{cor:sharing-density} & Necessary global continuation density & \leanname{cor:sharing-density} \\[2pt]
\cref{prop:metric-sharing} & Metric shared-support criterion & \leanname{prop:metric-sharing} \\[2pt]
\cref{cor:covering-exponents} & Sublinear raw covering exponents & \leanname{cor:covering-exponents} \\[2pt]
\cref{prop:serial-quadratic-cover} & Serial exceptional spans can require a quadratic cover & \leanname{prop:serial-quadratic-cover} \\[2pt]
\cref{thm:movement} & Theorem M: movement to cover distinct query types & \leanname{thm:movement} \\[2pt]
\cref{cor:movement-interval} & Queries spread over a column interval & \leanname{cor:movement-interval} \\[2pt]
\cref{prop:subset-calibration} & Exact helper and gap-area calibration & \leanname{prop:subset-calibration} \\[2pt]
\cref{cor:subset-movement} & The dimension limit precedes the radius limit & \leanname{cor:subset-movement} \\[2pt]
\cref{def:causal-group-typing} & Causal group typing & \leanname{def:causal-group-typing} \\[2pt]
\cref{thm:causal-group-routing} & Theorem Q: paid routing from causal groups & \leanname{thm:causal-group-routing} \\[2pt]
\cref{lem:attached-short-count}${}^{*}$ & Attached short-distance count & \leanname{lem:attached-short-count} \\*[2pt]
\cref{cor:causal-group-obstruction} & Narrow causal groups require an almost complete helper & \leanname{cor:causal-group-obstruction} \\
\end{longtable}
\endgroup

\noindent\begin{minipage}{\textwidth}
\paragraph{Verified declarations.}
The following declarations are audited by the generated \path{lean/PaperAxioms.lean}.
They depend at most on \leanname{propext}, \leanname{Classical.choice} and
\leanname{Quot.sound}. Earlier-paper declarations cited for reuse are verified there.\par
\noindent\leanfile{SourceCovers/FamilyAccounting.lean}:\par
{\small\raggedright\leanlist{SourceCovers.actualOwner,SourceCovers.ActualFamily,SourceCovers.HasSpanColoring,SourceCovers.mem_family_edges_iff,SourceCovers.constant_extension_isWalk}\par}
\noindent\leanfile{SourceCovers/FamilyAccounting.lean}:\par
{\small\raggedright\leanlist{SourceCovers.family_owner_inside,SourceCovers.family_hops_nested,SourceCovers.family_span_hop,SourceCovers.family_size_le_owner_arms,SourceCovers.two_mul_sum_family_size_le,SourceCovers.sum_family_variation_eq,SourceCovers.q2_accounting,SourceCovers.mem_boundaryLowers,SourceCovers.boundary_lowers_le_crossings,SourceCovers.boundary_lowers_le_family_sizes,SourceCovers.boundary_lowers_le}\par}
\noindent\leanfile{SourceCovers/ClippedShadows.lean}:\par
{\small\raggedright\leanlist{SourceCovers.ActualFamily.shadow,SourceCovers.ActualFamily.shadow_at_hop,SourceCovers.ActualFamily.shadow_exact,SourceCovers.shadowOrbit}\par}
\noindent\leanfile{SourceCovers/ClippedShadows.lean}:\par
{\small\raggedright\leanlist{SourceCovers.jointShadow_eq_on,SourceCovers.jointShadow_variation,SourceCovers.shadow_success_cons_iff,SourceCovers.shadow_word_structure,SourceCovers.shadow_word_variation,SourceCovers.variation_comp_le_crossings,SourceCovers.exists_sharp_patch}\par}
\noindent\leanfile{SourceCovers/ProtectedPhases.lean}:\par
{\small\raggedright\leanlist{SourceCovers.q2_caterpillar,SourceCovers.protected_survives,SourceCovers.admissiblePhases_card,SourceCovers.phase_payment_sum,SourceCovers.phase_average_exact,SourceCovers.protected_phase_average,SourceCovers.nonleaf_mass_le_total,SourceCovers.phase_path_length,SourceCovers.phase_success_length,SourceCovers.phase_shadowOrbit_exact,SourceCovers.phase_orbit_chains,SourceCovers.phase_orbit_card}\par}
\noindent\leanfile{SourceCovers/ExactSourceBank.lean}:\par
{\small\raggedright\leanlist{SourceCovers.exact_source_bank,SourceCovers.exactSourceBank_length,SourceCovers.exactSourceBank_bounded,SourceCovers.exactSourceBank_exact,SourceCovers.exactSourceBank_cover,SourceCovers.exactSourceBank_raw_root,SourceCovers.exactSourceBank_raw_values_card,SourceCovers.exactSourceBank_raw_selected_edge,SourceCovers.exactSourceBank_raw_variation}\par}
\noindent\leanfile{SourceCovers/IntervalCover.lean}:\par
{\small\raggedright\leanlist{SourceCovers.exists_cover_packing,SourceCovers.exists_subcover_le_weight}\par}
\noindent\leanfile{SourceCovers/NativeShortCount.lean}:\par
{\small\raggedright\leanlist{SourceCovers.attachedShortBits,SourceCovers.card_attachedShortBits_le_local,SourceCovers.card_attachedShortBits_le_profile,SourceCovers.card_attachedShortBits_le,SourceCovers.attachedShortBits_zero_radius,SourceCovers.attachedShortBits_zero_block}\par}
\end{minipage}\par

\section{Building and checking this draft}\label{sec:verification}

The paper sources and PDF use CC BY 4.0; the Lean library uses Apache 2.0.
\Cref{def:complete-families,lem:family-accounting,def:clipped-shadow,lem:shadow-variation,lem:protected-phase,thm:exact-source-bank,lem:interval-rounding,lem:attached-short-count}
are Lean-verified. Their cited
declarations depend only on \leanname{propext},
\leanname{Classical.choice} and \leanname{Quot.sound}. The remaining
numbered results have complete written proofs and are not formalized.
No additional axiom, admitted proof or native decision procedure is
used in the new library.

The package pins Lean and Mathlib to the repository's existing v4.34.1
versions. It requires the existing second-paper build adapter, whose
source directory points to the first paper's canonical Lean files.
It also requires the eighth paper's \path{boundedNesting} package for
actual families and the shared \path{lean-machine/} package for interval
covers, finite incidence sums, monotone patches, thin interval graphs,
protected phases, finite action orbits and uniform source banks.
Readers need the whole repository, with sibling folders in place.
No earlier mathematical source is copied or modified. From \path{lean/}, run
\begin{samepage}
\begin{verbatim}
lake exe cache get
lake build
lake env lean PaperAxioms.lean
\end{verbatim}
\end{samepage}
The cache command is needed only when the pinned dependency objects
are not already available. The axiom output must contain no axioms
beyond the three listed above. No heartbeat override is needed for
the new proofs.

From this paper's top-level directory, regenerate or check the statement
index and axiom commands with
\begin{verbatim}
python3 tools/statement_index.py
python3 tools/statement_index.py --check
\end{verbatim}
The script reads the section order from \path{paper/main.tex}, checks
stable labels, proof-status markers and the corresponding proof or
outline environments, and generates the index appendix,
\path{FORMALIZATION-INDEX.md}, and \path{lean/PaperAxioms.lean}.
It does not itself verify the mathematics or execute the axiom commands.
\path{FORMALIZATION.md} records completed work and the remaining plan.

Build the PDF from \path{paper/} with
\begin{verbatim}
latexmk -pdf -interaction=nonstopmode -halt-on-error main.tex
\end{verbatim}
Resolve undefined references and citations, and inspect the rendered
pages. Earlier-paper references are pinned in \path{paper/references.bib},
and their proof status is stated where used. Public paths are relative
to this repository. Run \path{tools/check_public_paths.py} from the
repository root as described in \path{CONTRIBUTING.md} before publishing.
That scan checks portability and local-path disclosure, not theorem
correctness.

\clearpage
\begingroup\small
\bibliographystyle{plainurl}
\bibliography{references}
\endgroup
\end{document}
